\documentclass[11pt]{article}
\usepackage[T1]{fontenc}
\usepackage[utf8]{inputenc}
\usepackage{lmodern,microtype}
\usepackage[margin=0.88in]{geometry}
\usepackage{amsmath,amssymb,amsthm,mathtools,mathrsfs}
\usepackage{array,longtable,booktabs,enumitem,xcolor}
\usepackage{hyperref,aliascnt}
\usepackage[nameinlink,capitalise,noabbrev]{cleveref}
\usepackage{fancyhdr}
\hypersetup{colorlinks=true,linkcolor=blue!35!black,citecolor=blue!35!black,urlcolor=blue!35!black,pdftitle={A proof blueprint for Auto.KoszAdjoint},pdfauthor={}}
\setlength{\parindent}{0pt}
\setlength{\parskip}{0.4em}
\setlength{\headheight}{14pt}
\setlength{\emergencystretch}{3em}
\pagestyle{fancy}\fancyhf{}\lhead{\small A proof blueprint for \texttt{Auto.KoszAdjoint}}\rhead{\small \thepage}
\newtheorem{theorem}{Theorem}[section]
\newaliascnt{proposition}{theorem}\newtheorem{proposition}[proposition]{Proposition}\aliascntresetthe{proposition}
\newaliascnt{lemma}{theorem}\newtheorem{lemma}[lemma]{Lemma}\aliascntresetthe{lemma}
\newaliascnt{corollary}{theorem}\newtheorem{corollary}[corollary]{Corollary}\aliascntresetthe{corollary}
\theoremstyle{definition}
\newaliascnt{definition}{theorem}\newtheorem{definition}[definition]{Definition}\aliascntresetthe{definition}
\newaliascnt{remark}{theorem}\newtheorem{remark}[remark]{Remark}\aliascntresetthe{remark}
\newcommand{\ee}[1]{\mathrm e\!\left(#1\right)}
\newcommand{\pid}[1]{\nolinkurl{#1}}
\newcommand{\one}{\mathbf 1}
\newcommand{\R}{\mathbb R}
\newcommand{\N}{\mathbb N}
\newcommand{\dd}{\,\mathrm d}
\newcommand{\Idx}{\{0,1,2,3\}}
\begin{document}
\begin{titlepage}
\vspace*{1.3cm}
{\Huge\bfseries A proof blueprint for\par \texttt{Auto.KoszAdjoint}\par}
\vspace{0.9cm}
{\Large The real monomial system $t,t^2,t^3$\par}
\vspace{0.5cm}
{\large Complete mathematical dependency development, including PET\par}
\vspace{1.1cm}
\textbf{Repository examined}\par
\texttt{roos-j/lean-dfr}, supplied \texttt{lean-dfr.task2.zip}.\par
\textbf{Pinned toolchain}\par
\texttt{leanprover/lean4:v4.33.0-rc1}.\par
\textbf{Pinned mathlib revision}\par
{\small\pid{79d0395a1825a6264ad5d269e35e60537518955e}}.\par
\textbf{Source}\par
Kosz--Mirek--Peluse--Wan--Wright, arXiv:2411.09478v4,\par
especially Sections 4--6.\par
\vfill
\textbf{Status of this document}\par
A mathematical proof and implementation specification, not a completed Lean formalization.
The supplied declarations and dependency manifest were inspected; this document has not been checked
by the Lean kernel. New declaration names in the implementation register are proposed names, not claims
about existing mathlib APIs.
\vspace{0.6cm}
\end{titlepage}
\setcounter{tocdepth}{2}\tableofcontents\clearpage
\section*{Scope and the exact remaining target}
\addcontentsline{toc}{section}{Scope and the exact remaining target}
The target is the proposition \pid{Auto.KoszAdjoint}, defined at lines 1803--1820 of
\pid{DFR/Auto/SmoothingIneq3D/Smoothing3D.lean} in the supplied archive.
Its three indices are represented in Lean by \pid{Fin 3}; the mathematical index $j\in\{1,2,3\}$
here corresponds to the Lean index $j-1$ and to \pid{expo (j-1) = j}.
No inverse theorem, major-arc theorem, polynomial-uniformity theorem, or adjoint estimate is assumed.
The existing improving estimate is also proved mathematically below, so that it is not an unexplained
analytic input. The subunit estimate is not needed for this target.

The dependency order is: elementary integration and Fourier identities; a real improving estimate;
signed van der Corput and a finite polynomial exhaustion argument; Fourier selection and degree lowering;
a noncircular major-arc induction; lowest-input and then every-input positive energy;
a compact high-frequency estimate; interpolation; support removal; and the exact extended-norm conclusion.
Statements introducing a conditional implication are proved before the induction supplying their premise.
The target is stated at the end, after all the facts used in its proof.

The inverse-theorem development reorganizes and expands the supplied \pid{patch_3_updated.tex}.
In particular, it retains the distinguished input through every polynomial operation and proves the
terminal cube identity, instead of appealing to an unspecified sequence of Cauchy--Schwarz applications.
The final adjoint argument uses a truncated high-frequency residual, rather than the older
Hahn--Banach structural decomposition. Thus Hahn--Banach separation, weak compactness of a
structured class, and a decomposition certificate are not additional obligations on this proof path.

Throughout, elementary facts about Lebesgue integration, change of variables, finite-dimensional
calculus, and finite polynomial algebra have their usual meanings. The precise nontrivial Fourier
identities for nonsmooth functions are proved below; they are not inferred from a theorem that only
identifies the Fourier transform of Schwartz functions. The implementation register states the
library interface and distinguishes checked-in declarations from newly required lemmas.


\section{Conventions and quantitative bookkeeping}\label{patch:conventions}
We use $\ee u=\exp(2\pi i u)$, $\langle f,g\rangle=\int f\overline g$, and
$\widehat f(\xi)=\int f(x)\ee{-x\cdot\xi}\,dx$.  Write $\mathcal C$ for complex conjugation.
For $N>0$,
\[
 A_N(f_1,f_2,f_3)(x)=\frac1N\int_0^N\prod_{i=1}^3f_i(x+t^ie_i)\,dt.
\]
and $\Lambda_N=\langle A_N(f_1,f_2,f_3),f_0\rangle$.
The notation $\|f\|_p$ denotes the Lebesgue $L^p$ norm.
Fix an even real smooth cutoff $\eta$ with
$\mathbf1_{[-1/4,1/4]}\leq\eta\leq\mathbf1_{(-1/2,1/2)}$.
The operator $P_L^{(j)}$ has multiplier $\eta(\xi_j/L)$.

Work first with Borel representatives, pointwise bounded by the stated bounds and zero off
the stated boxes.  Almost-everywhere versions follow at the end: for a null set $E$ and any
measurable translation $v(t)$,
\[
 \int\!\int\mathbf1_E(x+v(t))\,dx\,dt=0
\]
by translation invariance and Tonelli.  The same argument handles any finite number of shift
parameters.  Thus changing representatives does not change any correlation or cube integral.
All auxiliary phase functions can be chosen Borel after discarding the standard Fubini null
sets.

\subsection{A precise meaning for the power losses}
For an integer $c\geq1$, set
\[
 \ell_c(\delta)=c^{-1}\delta^c,\qquad u_c(\delta)=c\delta^{-c}
 \qquad(0<\delta\leq1).
\]
An assertion below of the form $X\geq\delta^{O(1)}S$ means that the proof returns an integer
$c$, depending only on the fixed degrees, support/coefficient budgets, and any fixed cube order,
with $X\geq\ell_c(\delta)S$.  Similarly $X\leq\delta^{-O(1)}S$ returns $c$ with
$X\leq u_c(\delta)S$.  It never means that the exponent may depend on a function, a selected
section, $N$, or the numerical value of $\delta$.

For implementation, retain the returned integers rather than using asymptotic notation.
The following rules suffice to compose every budget in the proof.
\begin{lemma}[Budget arithmetic and popularity]\label{patch:budgets}
Finite sums, products, quotients with a positive lower bound, and fixed natural powers of
power-controlled positive quantities are power-controlled.  A fixed positive root of a
power lower bound is again a power lower bound.  If a finite measure space has measure at most
$M$, $0\leq F\leq B$, and $\int F\geq a>0$, then
\[
 E=\{F\geq a/(2M)\}\quad\hbox{satisfies}\quad \mu(E)\geq a/(2B).
\]
Removing a set of measure at most $a/(2B)$ loses at most half the lower bound.  A partition
into $K$ measurable cells has a cell carrying at least $1/K$ of a nonnegative total mass.
\end{lemma}
\begin{proof}
For upper bounds choose, for example, an integer at least both the new coefficient and the
new exponent.  For a product of $u_a$ and $u_b$ it suffices to take
$c\geq\max(ab,a+b)$; for a product of $\ell_a$ and $\ell_b$ the same choice works.
For $r$-th powers take $c\geq\max(a^r,ra)$.  Bound $X/Y$ using an upper bound for $X$
and a positive lower bound for $Y$, or the reverse for a lower bound.  Round all constants
upward by the Archimedean property.  Since $0<\delta\leq1$, increasing a lower-bound exponent
weakens it and increasing an upper-bound exponent enlarges it.  These observations also
handle roots and rational fixed exponents.
For popularity, the integral on $E^c$ is at most $a/2$, whereas the integral on $E$ is at most
$B\mu(E)$.  The remaining assertions follow from integral monotonicity and a finite sum.
\end{proof}

Whenever a proof chooses a finer parameter, it uses these explicit recipes:
\[
 \text{discarded measure}\leq\frac{\text{retained mass}}{4\,\text{integrand bound}},
 \qquad
 \text{phase error}\leq\frac{\text{retained mass}}{4\,\text{integration-volume bound}}.
\]
For a phase perturbation use $|\ee u-\ee v|\leq2\pi|u-v|$.
A bounded interval of frequencies of length $2M$ can be partitioned into at most
$1+\lceil2M/\tau\rceil$ cells of width $\tau$.  All such counts below are power-controlled.
The iteration counts are bounded independently of $\delta$ before these recipes are applied.

\paragraph{Scale convention in the inverse module.}
For $1\leq d_1<\cdots<d_k\leq3$, write $D_r=\sum_{i\leq r}d_i$, $D_0=0$ and $D=D_k$.
A budgeted support box has sides $[-u_c(\delta)N^{d_i},u_c(\delta)N^{d_i}]$.
A polynomial $P_i$ is admissible with budget $c$ when
\[
 \deg P_i=d_i,\qquad
 \ell_c(\delta)\leq|\operatorname{lc}P_i|\leq u_c(\delta),\qquad
 |[t^r]P_i|\leq u_c(\delta)N^{d_i-r}\quad(0\leq r<d_i).
\]
These are the paper's admissibility and enlarged-support hypotheses, with a single integer
budget enlarged as needed.  Setting $t=Nu$ and $x_i=N^{d_i}y_i$ reduces all coefficient,
volume, and sublevel estimates to bounded parameter intervals with power-controlled
constants.  This change has determinant $N^D$; the powers of $N$ in subsequent statements
are therefore exact, not part of the budget.



\paragraph{All uses of power bounds are quantified.}
In the inverse module, ``power-large'' means a lower bound $\ell_c(\delta)$ times the
stated physical normalization, for a specified input budget $c$. ``Power-controlled'' means
between such a lower bound and $u_c(\delta)$ times that normalization. Each theorem is
quantified over every input budget and returns an integer output budget. Every choice made
below is one of the arithmetic, threshold, or finite-partition operations in
\cref{patch:budgets}, except the finite polynomial recursion, whose independent uniform
bound is proved explicitly. Thus no choice may depend on a selected function value or on
an unbounded family of parameter sections. The notation $\delta^{O(1)}$ and
$\delta^{-O(1)}$ is only a display abbreviation for these returned integers, not an
additional hypothesis or an unspecified uniformity principle.

For clarity, an application stated with a ``power lower bound'' has the logical form:
for every input budget $c$ there is an output budget $C$ such that the hypotheses with
$\ell_c,u_c$ imply the stated conclusion with $\ell_C,u_C$. This is also the meaning for
all correlations and section measures in the intermediate lemmas. At each application
feed it the common lower threshold from the preceding popularity step, not the actual
value on a particular section. Enlarging a cutoff radius to compare energies always uses
at least a factor of two; monotonicity of $\eta$ is never silently invoked.


\section{Analytic facts used by the proof}\label{sec:analytic-foundations}
All measures in this section are Lebesgue measures. The space $\R^0$ is a singleton with
measure one. A box in $\R^n$ is a product of bounded intervals; its measure is the product
of their lengths. All equalities of $L^p$ functions are understood almost everywhere unless
pointwise equality is expressly stated. Write
\[
 k(u)=\int_{\R}\eta(\xi)\ee{u\xi}\,d\xi,
 \qquad K_R(u)=Rk(Ru)\quad(R>0).
\]
The cutoff is the one fixed in the conventions. It is available in the supplied source as
\pid{Auto.eta}, constructed from \pid{ContDiffBump} with inner radius $1/4$ and outer radius $1/2$.
No monotonicity of $\eta$ as a function of $|\xi|$ is assumed.

\subsection{Integration, approximation, and norm inequalities}
\begin{lemma}[The elementary integral inequalities in the required form]\label{new:integral-inequalities}
Let $1\leq p<\infty$. Translations preserve the $L^p$ norm. If $1/p=\sum_{i=1}^r1/p_i$
with $p_i\in[1,\infty]$, then
\[
 \left\|\prod_{i=1}^r f_i\right\|_p\leq\prod_{i=1}^r\|f_i\|_{p_i}.
\]
For jointly measurable $F$ with $\int\|F(u,\cdot)\|_p\,d\nu(u)<\infty$,
\[
 \left\|\int F(u,\cdot)\,d\nu(u)\right\|_p
 \leq\int\|F(u,\cdot)\|_p\,d\nu(u).
\]
Consequently, if $a\in L^1(\R)$ and $f\in L^p(\R^n)$, their one-coordinate convolution satisfies
\[
 \left\|\int a(u)f(\cdot-ue_j)\,du\right\|_p\leq\|a\|_1\|f\|_p.
\]
The analogous $p=\infty$ statements hold with essential suprema.
\end{lemma}
\begin{proof}
Translation invariance is the substitution $y=x+v$ in the integral of $|f|^p$.
For the product inequality apply H\"older to $\prod |f_i|^p$ with exponents $p_i/p$;
factors with exponent $\infty$ are removed by their essential supremum. Induction gives the
finite-factor version from the two-factor version. For $p=1$, the integral inequality is
Tonelli and $|\int F|\leq\int|F|$. For $1<p<\infty$, first prove it on finite sums:
H\"older applied to $|\sum F_i|^{p-1}$ gives the usual triangle inequality after division by
$\|\sum F_i\|_p^{p-1}$, with the zero case separated. Simple approximation in the parameter,
followed by Fatou, gives the displayed integral inequality. Equivalently, one may use the
Bochner integral inequality in the Banach space $L^p$; strong measurability follows by simple
approximation of the joint measurable function. The convolution conclusion follows because
$\|a(u)f(\cdot-ue_j)\|_p=|a(u)|\|f\|_p$. For $p=\infty$ use the pointwise integral bound
outside the joint null set supplied by Tonelli.
\end{proof}

\begin{lemma}[Approximation and translation continuity]\label{new:translation-continuity}
If $1\leq p<\infty$ and $f\in L^p(\R^n)$, bounded finite-valued functions supported in bounded
boxes approximate $f$ in $L^p$. Moreover,
\[
 \|f(\cdot-h)-f\|_p\longrightarrow0\quad\text{as }h\longrightarrow0.
\]
\end{lemma}
\begin{proof}
First replace $f$ by $f\one_{[-M,M]^n}\one_{\{|f|\leq M\}}$. The $p$-th power of the
error tends to zero by dominated convergence, with majorant $|f|^p$. On this bounded box,
round the real and imaginary parts to multiples of $1/q$ within their bounded ranges.
The rounding error is at most $\sqrt2/q$ pointwise, so its $L^p$ norm is at most
$\sqrt2\,|[-M,M]^n|^{1/p}/q$. These are the asserted approximants.

For the indicator of a finite-measure measurable set $E$, fix $\epsilon>0$. Lebesgue
regularity gives compact $K\subseteq E$ with $|E\setminus K|<\epsilon$, and a bounded open
$O\supseteq K$ with $|O\setminus K|<\epsilon$. Compactness gives a neighborhood of zero
such that $K+h\subseteq O$ there: cover $K$ by finitely many balls whose doubled radii
remain inside $O$ and take the minimum radius. Since $|K+h|=|K|$,
\[
 |K\mathbin\triangle(K+h)|=2|(K+h)\setminus K|\leq2|O\setminus K|<2\epsilon.
\]
The symmetric-difference triangle inequality therefore gives
$|E\mathbin\triangle(E+h)|<4\epsilon$. This equals
$\|\one_E(\cdot-h)-\one_E\|_p^p$. Finite linear combinations inherit translation continuity
by the triangle inequality. For general $f$, choose a simple approximant $s$ and use
\[
 \|f(\cdot-h)-f\|_p\leq2\|f-s\|_p+\|s(\cdot-h)-s\|_p.
\]
First make the first term small and then let $h$ tend to zero.
\end{proof}

\begin{lemma}[Duality in the form used here]\label{new:duality}
For $1<p<\infty$, $p'=p/(p-1)$, and a measurable $F$, a bound
\[
 \left|\int Fh\right|\leq B\|h\|_{p'}
\]
for all bounded, compactly supported measurable $h$ for which the pairing is defined implies
$F\in L^p$ and $\|F\|_p\leq B$. It is enough to test finite-valued such $h$ when $F$ is locally
integrable. The same norm bound, with the same class of tests, holds whenever $F\in L^p$.
\end{lemma}
\begin{proof}
The upper bound when $F\in L^p$ is H\"older. For the converse put
$E_M=[-M,M]^n\cap\{M^{-1}\leq|F|\leq M\}$ and
$h_M=\one_{E_M}\overline F|F|^{p-2}$. It is bounded with bounded support, and
\[
 \int Fh_M=\int_{E_M}|F|^p,
 \qquad \|h_M\|_{p'}=\left(\int_{E_M}|F|^p\right)^{1/p'}.
\]
If this integral is positive, division gives its $1/p$ power at most $B$; if it is zero the
same conclusion holds. The increasing sets $E_M$ exhaust $\{0<|F|<\infty\}$, so monotone
convergence proves the result. For simple tests only, approximate each $h_M$ uniformly on
its support by finite-valued functions. Local integrability of $F$ justifies passage in the
pairing, and uniform approximation on a finite-measure support justifies passage in the
$L^{p'}$ norm. The functions arising before the extension steps below are bounded and in
$L^1$, so this local integrability hypothesis is always satisfied.
\end{proof}

\begin{lemma}[Null sets, sections, and measurable unit phases]\label{new:measurability}
All finite products, translates, averaged translates, and cube averages of the bounded Borel
functions used below are jointly measurable. For a measurable scalar function $b$, the function
$\overline b/|b|$ on $\{b\ne0\}$ and zero elsewhere is measurable, bounded by one, and
its product with $b$ equals $|b|$. Replacing any input by an almost-everywhere equal input
does not change the averaged functions almost everywhere or any of their integrals.
\end{lemma}
\begin{proof}
The translation maps used below are continuous polynomial maps in their real variables;
composition with Borel functions and finite products are Borel. Parameter integrals are
measurable by the measurable-integral theorem, applied first to nonnegative real and
imaginary parts and then to integrable complex functions. Absolute integrability follows
from the explicitly stated bounded-support envelopes or from
\cref{new:integral-inequalities}. For the phase, division is continuous away from zero and
the two pieces are measurable. Finally, for a null set $E$ and any measurable translation
$v(t)$, Tonelli and translation invariance give
\[
 \int\!\int\one_E(x+v(t))\,dx\,dt=\int|E|\,dt=0.
\]
The same calculation applies to finitely many parameter variables. Lebesgue measurable
functions admit Borel versions; changing to these versions is therefore harmless. Every
sectional identity below is used on the full-measure set furnished by Fubini, not on an
unjustified collection of all exceptional sections.
\end{proof}

\subsection{Schwartz kernels and integral Plancherel without a missing compatibility lemma}
\begin{lemma}[Kernel facts]\label{new:kernel-facts}
The function $k$ is real, even, and Schwartz, $\widehat k=\eta$, and $\int k=1$.
For every integer $M\geq0$, the quantities
\[
 B_0=\|k\|_1,\qquad B_1=\|k'\|_1,\qquad
 B_M^{\rm tail}=\int |u|^M|k(u)|\,du
\]
are finite. For $R,a>0$,
\[
 \|K_R\|_1=B_0,\qquad \|K_R'\|_1=RB_1,
 \qquad \int_{|u|\geq a}|K_R(u)|\,du\leq B_M^{\rm tail}(Ra)^{-M}.
\]
Also, if $|f|\leq1$, then $P_R^{(j)}f=K_R*_{j}f$ satisfies
\[
 \|P_R^{(j)}f\|_\infty\leq B_0,
 \qquad |P_R^{(j)}f(x+ve_j)-P_R^{(j)}f(x)|\leq RB_1|v|.
\]
\end{lemma}
\begin{proof}
A smooth compactly supported function is Schwartz. The Fourier transform and inverse
Fourier transform preserve Schwartz space, and on that space their compositions are the
identity. These are the Schwartz Fourier equivalence and inversion already used by
\pid{Auto.fourier_etaKer}; no nonsmooth Plancherel identification enters here.
Alternatively, decay of $k$ and each of its derivatives follows by integrating by parts
arbitrarily many times in its defining integral; the boundary terms vanish because $\eta$
has compact support and is smooth. The imaginary part is zero by oddness of the sine term,
and changing $u$ to $-u$ proves evenness. Evaluating $\widehat k=\eta$ at zero gives
$\int k=\eta(0)=1$. Schwartz decay proves all displayed integrability statements.
The substitutions $v=Ru$ give the first two scaled norms. On $|v|\geq Ra$ use
$1\leq |v|^M/(Ra)^M$ to obtain the tail inequality. The supremum bound follows from the
integral triangle inequality. For the Lipschitz bound use
\[
 \int|K_R(u+v)-K_R(u)|\,du
 \leq\int\!\int_0^{|v|}|K_R'(u+\operatorname{sgn}(v)t)|\,dt\,du
 =|v|RB_1.
\]
Translation in the convolution variable proves the claimed bound for every bounded Borel
section. This argument does not require differentiability of $f$.
\end{proof}

\begin{lemma}[An $L^2$ approximate identity]\label{new:approximate-identity}
Let $a\in L^1(\R^n)$ with $\int a=1$ and put $a_R(x)=R^na(Rx)$. For every $f\in L^2$,
\[
 \|a_R*f-f\|_2\longrightarrow0\quad(R\longrightarrow\infty).
\]
\end{lemma}
\begin{proof}
Substitute $u=Ry$ in the convolution. The integral inequality gives
\[
 \|a_R*f-f\|_2\leq\int|a(u)|\,\|f(\cdot-u/R)-f\|_2\,du.
\]
For each $u$ the norm tends to zero by \cref{new:translation-continuity}, and it is at most
$2\|f\|_2$. Dominated convergence with the integrable majorant
$2\|f\|_2|a(u)|$ proves the assertion.
\end{proof}

\begin{theorem}[Plancherel for the integral Fourier transform]\label{new:integral-plancherel}
If $f\in L^1(\R^n)\cap L^2(\R^n)$, define its Fourier transform by the ordinary integral.
Then $\widehat f\in L^2$ and
\[
 \int|\widehat f(\xi)|^2\,d\xi=\int|f(x)|^2\,dx.
\]
For $f,g\in L^1\cap L^2$,
\[
 \int f\overline g=\int\widehat f\,\overline{\widehat g}.
\]
\end{theorem}
\begin{proof}
Put $\chi(\xi)=\prod_{i=1}^n\eta(\xi_i)$ and
$a(x)=\prod_{i=1}^nk(x_i)$. Thus $a\in L^1$, $\int a=1$, and $a$ is the inverse Fourier
transform of $\chi$, by repeated Fubini. For $R>0$ set $\chi_R(\xi)=\chi(\xi/R)$ and
$a_R(x)=R^na(Rx)$. Expanding the two Fourier integrals gives
\[
 \int\chi_R(\xi)|\widehat f(\xi)|^2\,d\xi
 =\int f(x)\overline{(a_R*f)(x)}\,dx.
\]
Indeed, the expanded absolute integral is bounded by
$\|\chi_R\|_1\|f\|_1^2<\infty$, so Fubini is applicable. The inner frequency integral is
$a_R(y-x)=a_R(x-y)$, and it is real; this gives exactly the displayed pairing.
By \cref{new:approximate-identity} and Cauchy--Schwarz its right side tends to $\|f\|_2^2$.
Take $R=n+1$ with $n$ a natural number. The left integrands are nonnegative and converge
pointwise to $|\widehat f|^2$, because
$\chi_R\to1$. Fatou first gives $\int|\widehat f|^2\leq\|f\|_2^2$, proving its integrability.
Now dominated convergence, with majorant $|\widehat f|^2$, proves equality. Notice that
monotonicity in $R$ was not used. Applying this norm identity to $f+g$, $f-g$, $f+ig$,
and $f-ig$, and expanding the four squares, proves the pairing identity by polarization.
All these sums remain in $L^1\cap L^2$.
\end{proof}

\begin{lemma}[Multipliers, directional energies, and slices]\label{new:multiplier-facts}
For $f\in L^1\cap L^2$ and $R>0$,
\[
 \widehat{P_R^{(j)}f}(\xi)=\eta(\xi_j/R)\widehat f(\xi),
 \qquad \|P_R^{(j)}f\|_2\leq\|f\|_2.
\]
The operator $P_R^{(j)}$ is self-adjoint, and
\[
 \langle f,P_R^{(j)}f\rangle
 =\int\eta(\xi_j/R)|\widehat f(\xi)|^2\,d\xi\geq0.
\]
If $L\leq R/2$, then $P_L^{(j)}(I-P_R^{(j)})f=0$ in $L^2$. If $L'\geq2L$, then
\[
 \langle f,P_{L'}^{(j)}f\rangle\geq\langle f,P_L^{(j)}f\rangle.
\]
For a coordinate $i\ne j$, writing $f_v$ for the section at $x_i=v$ gives
\[
 \langle f,P_R^{(j)}f\rangle=\int\langle f_v,P_R^{(j)}f_v\rangle\,dv.
\]
For bounded compactly supported $f$, each sectional energy is nonnegative and the last formula
holds with the sectionwise integral transform, on the Fubini full-measure set.
\end{lemma}
\begin{proof}
The absolute double integral for Fourier transformation of $K_R*_{j}f$ is at most
$\|K_R\|_1\|f\|_1$. Fubini and the translation $x=y+ue_j$ give the multiplier
$\widehat{K_R}(\xi_j)=\eta(\xi_j/R)$. The convolution belongs to $L^1\cap L^2$ by
\cref{new:integral-inequalities}. Apply \cref{new:integral-plancherel}; the bounds
$0\leq\eta\leq1$ give the contraction and energy identities. The pairing identity also
proves self-adjointness. On the support of $\eta(\xi_j/L)$, one has
$|\xi_j|\leq L/2\leq R/4$, so $\eta(\xi_j/R)=1$. Thus the annihilation multiplier is
zero; integral Plancherel makes its physical function zero in $L^2$.
Likewise the cutoff at $L'$ equals one on the support of the cutoff at $L$ when $L'\geq2L$;
this proves the energy comparison without radial monotonicity.

For slicing, expand the physical convolution pairing. The absolute integral is bounded by
$\|K_R\|_1\|f\|_2^2$, using Cauchy--Schwarz and translation invariance. Fubini therefore
allows the $v$ integral to be outermost and gives the displayed formula. The sectional
nonnegativity follows from the already proved energy identity in one lower dimension.
Translation or conjugation of $f$ preserves this energy, since the Fourier modulus is
unchanged up to reflection and the cutoff is even. The same arguments establish all these
identities for the $L^2$ extensions: approximate by $L^1\cap L^2$ functions and use the
contraction and Cauchy--Schwarz. The inverse argument below itself only needs $L^1\cap L^2$.
\end{proof}

\begin{lemma}[The averaged-translation multiplier]\label{new:average-multiplier}
Let $v(t)$ be a measurable bounded vector on a finite probability space and $|b(t)|\leq1$.
For $f\in L^1\cap L^2$ put $Tf(x)=\mathbb E_t b(t)f(x-v(t))$. Then $Tf\in L^1\cap L^2$,
\[
 \widehat{Tf}(\xi)=\left(\mathbb E_t b(t)\ee{-\xi\cdot v(t)}\right)\widehat f(\xi),
\]
and its squared $L^2$ norm is the integral of the squared modulus of the displayed multiplier
against $|\widehat f|^2$.
\end{lemma}
\begin{proof}
The $L^1$ and $L^2$ bounds are at most the corresponding norms of $f$, by
\cref{new:integral-inequalities}. The absolute joint Fourier integral is at most $\|f\|_1$.
Fubini and translation give the formula; \cref{new:integral-plancherel} gives the norm identity.
\end{proof}

\subsection{An interpolation lemma with proof}
\begin{lemma}[Three lines in the form used below]\label{new:three-lines}
Suppose $F$ is continuous and bounded on $0\leq\operatorname{Re}z\leq1$, holomorphic in
its interior, and $|F(it)|\leq M_0$, $|F(1+it)|\leq M_1$, with $M_0,M_1>0$.
For $0<\theta<1$,
\[
 |F(\theta)|\leq M_0^{1-\theta}M_1^\theta.
\]
\end{lemma}
\begin{proof}
Divide $F(z)$ by $\exp((1-z)\log M_0+z\log M_1)$, obtaining a bounded holomorphic function
$G$ with boundary modulus at most one. For $\epsilon>0$ consider
$G(z)\exp(\epsilon(z^2-1))$ on the rectangle $0\leq\operatorname{Re}z\leq1$,
$|\operatorname{Im}z|\leq T$. On the vertical sides its modulus is at most one. On the
horizontal sides it is at most $\sup|G|\exp(-\epsilon T^2)$, hence at most one for sufficiently
large $T$. The maximum-modulus theorem gives the same bound inside, so
$|G(\theta)|\leq\exp(\epsilon(1-\theta^2))$. Let $\epsilon$ decrease to zero.
\end{proof}

\begin{lemma}[Multilinear interpolation on a fixed support domain]\label{new:interpolation}
Let $T$ be a complex multilinear operator on bounded finite-valued functions of bounded support,
with $r$ inputs. The support restrictions may be fixed measurable sets, provided they are unchanged
when nonzero values of a simple function are replaced by other nonzero values.
Assume, for $a=0,1$, the estimates
\[
 \|T(f_1,\ldots,f_r)\|_{q_a}\leq M_a\prod_i\|f_i\|_{p_{i,a}},
\]
where $1<q_a<\infty$, $p_{i,a}\in[1,\infty]$, and $M_a>0$.
Define $1/p_i=(1-\theta)/p_{i,0}+\theta/p_{i,1}$ and
$1/q=(1-\theta)/q_0+\theta/q_1$, with $0<\theta<1$ and all $p_i<\infty$.
Then
\[
 \|T(f_1,\ldots,f_r)\|_q\leq M_0^{1-\theta}M_1^\theta\prod_i\|f_i\|_{p_i}.
\]
The same result applies when some inputs are conjugate-linear, by conjugating those inputs
first. It also applies to their norm-continuous extensions.
\end{lemma}
\begin{proof}
The zero-norm cases vanish by multilinearity. By homogeneity suppose $\|f_i\|_{p_i}=1$.
For a simple test function $h$ with $\|h\|_{q'}=1$, use the bilinear pairing $\int Th$, without
complex conjugation, and define on the nonzero values of each input
\[
 f_{i,z}=\frac{f_i}{|f_i|}\exp\!\left(p_i\left(\frac{1-z}{p_{i,0}}+
                   \frac{z}{p_{i,1}}\right)\log|f_i|\right).
\]
Set it equal to zero where $f_i=0$, and use $1/\infty=0$. Define $h_z$ in the same way using
the pair $q_0',q_1'$ and the exponent $q'$. The identity
$1/q'=(1-\theta)/q_0'+\theta/q_1'$ follows by subtracting the output identity from one.
At $z=\theta$ these are the original functions. At $z=a+it$, each input has $L^{p_{i,a}}$
norm one; when that exponent is infinite its modulus is its support indicator. Also
$\|h_{a+it}\|_{q_a'}=1$. Every function has its original fixed support.

Expand the simple functions on their disjoint level sets. The scalar
$F(z)=\int T(f_{1,z},\ldots,f_{r,z})h_z$ is a finite sum of exponentials of affine functions
of $z$, whose coefficients are the fixed pairings of indicator inputs and indicator tests.
These pairings are finite by the assumed endpoint bounds. Hence $F$ is entire and bounded
on the closed strip. H\"older and the two endpoint estimates bound its boundary values by
$M_0,M_1$. Apply \cref{new:three-lines}. Taking the supremum over simple tests and using
\cref{new:duality} gives the result. For general bounded supported inputs, simple approximation
and the existing crude multilinear bound justify passage to the limit in the applications
below; the precise extensions are carried out where they are needed.
\end{proof}


\section{A real improving estimate, proved before it is used}\label{new:improving-section}
This section proves the particular improving estimate needed at each adjoint index.
It also explains why no integer mean-value theorem is involved. The conclusion is already
available, in a more symmetric form, as \pid{Auto.strong_real_improving'} and
\pid{Auto.kosz53_internal} in the supplied Lean file. The proof here makes the mathematical
dependency explicit instead of treating either result as an imported analytic assertion.

\subsection{A six-dimensional incidence estimate}
Work at scale one in this subsection. Write $\R^6=\R\times\R^2\times\R^3$ and put
\[
 \Gamma_1(t)=(t;0,0;0,0,0),\qquad
 \Gamma_2(t)=(0;t,t^2;0,0,0),\qquad
 \Gamma_3(t)=(0;0,0;t,t^2,t^3),\qquad \Gamma_0=0.
\]
For finite-measure Borel sets $E_0,E_1,E_2,E_3\subseteq\R^6$ define the incidence
\[
 K=\int\one_{E_0}(x)\int_0^1\prod_{i=1}^3\one_{E_i}(x-\Gamma_i(t))\,dt\,dx.
\]
By translation and Tonelli, $0\leq K\leq\min_i|E_i|$. When $K>0$, write
$\alpha_i=K/|E_i|\in(0,1]$.

\begin{lemma}[Finite refinements and their pointwise incidence bounds]\label{new:refinements}
There exist nested Borel sets $E_i^r\subseteq E_i$, $0\leq r\leq60$, with $E_i^0=E_i$,
whose joint incidence after stage $r$ is at least $2^{-4r}K>0$.
For $x\in E_a^r$, the set of parameters taking $x$ to the required sets at the other labels
has measure at least $2^{-240}\alpha_a$. More precisely, for $b\ne a$, the translated point
$x+\Gamma_a(t)-\Gamma_b(t)$ lies in $E_b^r$ if $b<a$ and in $E_b^{r-1}$ if $b>a$,
simultaneously for all $b\ne a$, on such a parameter set.
\end{lemma}
\begin{proof}
At stage $r$, update labels $a=0,1,2,3$ in that order. Keep all already updated sets at level
$r$, and the other sets at level $r-1$. For label $a$, let $H_a^r(x)$ be the integral in
$t\in[0,1]$ of their incidence indicators after translating the point at label $a$ to $x$:
\[
 H_a^r(x)=\int_0^1\prod_{b<a}\one_{E_b^r}(x+\Gamma_a(t)-\Gamma_b(t))
                       \prod_{b>a}\one_{E_b^{r-1}}(x+\Gamma_a(t)-\Gamma_b(t))\,dt.
\]
Define
\[
 E_a^r=E_a^{r-1}\cap\{H_a^r\geq2^{-(4(r-1)+a+1)}\alpha_a\}.
\]
The functions and sets are measurable by \cref{new:measurability}. Before this deletion the
incidence is at least $2^{-(4(r-1)+a)}K$, by induction on the lexicographic pair $(r,a)$.
The deleted contribution is at most
$2^{-(4(r-1)+a+1)}\alpha_a|E_a|=2^{-(4(r-1)+a+1)}K$.
Thus at least half remains. The adjoint change of variables identifying the incidence at
label $a$ with the preceding joint incidence is exactly $y=x+\Gamma_a(t)$, of determinant one.
This proves the induction and the pointwise bounds. For $r\leq60$ the displayed threshold
is at least $2^{-240}\alpha_a$.
\end{proof}

\begin{lemma}[The two flows actually needed]\label{new:flows}
For each terminal label $b\in\{1,3\}$, there is a Borel parameter tower
$\mathcal P\subseteq[0,1]^6$ and a polynomial map $\Phi:\R^6\to\R^6$ whose image on
$\mathcal P$ lies in $E_b$. Each successive tower fiber has measure at least
$2^{-240}\alpha_a$, where $a$ is the label being left. In the order of the moment blocks,
\[
 |\det D\Phi|=12|u-v|\,|a_1-a_2|\,|a_1-a_3|\,|a_2-a_3|.
\]
Here $u,v$ are the two parameters of the degree-two block and $a_1,a_2,a_3$ the three
parameters of the degree-three block. On the set where this determinant is nonzero,
$\Phi$ is the union of two injective restrictions.
\end{lemma}
\begin{proof}
Use the following fixed sequences of labels, respectively for terminal labels $3$ and $1$:
\[
 (1,0,2,0,3,0,3),\qquad (3,0,3,0,2,0,1).
\]
Choose a point $x_0$ in the level-60 refinement of the initial label; it is nonempty because
the joint incidence there is positive. At a move from label $a$ at level $r$ to label $b$,
use the parameter fiber from \cref{new:refinements}, and put
\[
 x_{\ell}=x_{\ell-1}+\Gamma_a(t_\ell)-\Gamma_b(t_\ell).
\]
The new level is $r$ if $b<a$, and $r-1$ if $b>a$. Each of these two sequences has three
increases of label, so no level below 57 is used. Define the tower recursively by these
fiber conditions. Each condition is Borel in the previous parameters, giving a Borel tower;
the fiber lower bound is exactly the one at the label being left. The endpoint map is
\[
 \Phi(t_1,\ldots,t_6)=x_0+\sum_{\ell=1}^6
       (\Gamma_{w_{\ell-1}}(t_\ell)-\Gamma_{w_\ell}(t_\ell)).
\]
Its image lies in the required terminal refinement.

For the first sequence the independent moment blocks use parameters
$(t_1)$, $(t_2,t_3)$, $(t_4,t_5,t_6)$. For the second they use
$(t_1,t_2,t_3)$, $(t_4,t_5)$, $(t_6)$, in degrees three, two, one.
Every move touches label zero, so there are no cross-block terms. Up to a common sign in
each block, the three maps are
\[
 u\longmapsto u,\qquad
 (u,v)\longmapsto(u-v,u^2-v^2),
\]
\[
 (a_1,a_2,a_3)\longmapsto
 (a_1-a_2+a_3,a_1^2-a_2^2+a_3^2,a_1^3-a_2^3+a_3^3).
\]
Their derivative determinants have absolute values $1$, $2|u-v|$, and
$6|a_1-a_2||a_1-a_3||a_2-a_3|$, respectively. The last formula follows by factoring
$1,2,3$ from the derivative rows and expanding the three-column Vandermonde determinant.
Multiplication gives the claimed formula.

To verify multiplicity, the first block determines its parameter. In the second, if
$s_1=u-v\ne0$ and $s_2=u^2-v^2$, then $u+v=s_2/s_1$, determining both parameters.
For the third put $U=a_1+a_3$, $V=a_1a_3$, and let $s_1,s_2,s_3$ be its three outputs.
Direct expansion gives
\[
 a_2=U-s_1,\qquad V=s_1U-\frac{s_1^2+s_2}{2},
\]
\[
 s_3=\frac32(s_2-s_1^2)U+s_1^3,
 \qquad s_2-s_1^2=2(a_1-a_2)(a_2-a_3).
\]
At a noncritical point the last expression is nonzero. Thus $U$, then $V$, then $a_2$
are determined. The distinct numbers $a_1,a_3$ are the two roots of $X^2-UX+V$, so there
are at most two orderings. Restricting to $a_1<a_3$ or $a_1>a_3$ gives two Borel injective
pieces. This proves the assertion without a general polynomial multiplicity theorem.
\end{proof}

\begin{lemma}[Integrated incidence consequence]\label{new:incidence-consequence}
There is an absolute $c_*>0$ such that, for $b\in\{1,3\}$,
\[
 |E_b|\geq c_*(\alpha_0\alpha_1\alpha_2\alpha_3)^{10}.
\]
Consequently, for any $j\in\{1,2,3\}\setminus\{b\}$ and the remaining index $c$,
\[
 K\leq C_*|E_0|^{1/6}|E_j|^{1/2}|E_b|^{11/60}|E_c|^{1/6}.
\]
\end{lemma}
\begin{proof}
Let $\gamma=2^{-240}\alpha_0\alpha_1\alpha_2\alpha_3$. Every fiber in the preceding
lemma has measure at least $\gamma$, since the product of the four numbers $\alpha_i\leq1$
is at most each factor. When selecting a new parameter in a moment block, remove the open
intervals of radius $\gamma/8$ around the at most two previously selected parameters of
that block. Their total length is at most $\gamma/2$, so each pruned fiber has measure at
least $\gamma/2$. The conditions defining these removals are Borel, so the pruned set is
still a measurable tower. All four Vandermonde factors on it have modulus at least $\gamma/8$.
Successive integration over its six fibers gives
\[
 \int_{\mathcal P}|\det D\Phi|\geq
 12(\gamma/8)^4(\gamma/2)^6.
\]
Use the injective change-of-variables inequality on each of the two pieces in
\cref{new:flows}: for a $C^1$ map injective on a measurable set $S$,
$\int_S|\det D\Phi|\leq|\Phi(S)|$, with outer measure if necessary.
Each image is contained in $E_b$, so the whole integral is at most $2|E_b|$.
This is the usual Jacobian change-of-variables theorem; the precise one-sided form is
available as \pid{MeasureTheory.lintegral_abs_det_fderiv_le_addHaar_image}, used in the
supplied file. The critical set contributes zero. Thus one may take
$c_*=6\,8^{-4}2^{-6}2^{-2400}$.

Since $\alpha_j\leq1$, replacing its exponent 10 by 30 weakens the resulting lower bound:
\[
 |E_b|\geq c_*\alpha_0^{10}\alpha_j^{30}\alpha_b^{10}\alpha_c^{10}.
\]
Substitute $\alpha_i=K/|E_i|$, multiply by the denominator, and take the sixtieth root.
This gives the second assertion with $C_*=c_*^{-1/60}$. If $K=0$ the conclusion is immediate;
if a set has measure zero then $K=0$, as was noted before the refinements.
\end{proof}

\subsection{Descent and the strong endpoint, including the summation}
For the sign convention of this subsection put
\[
 T(f_1,f_2,f_3)(x)=\int_0^1\prod_{i=1}^3f_i(x-t^ie_i)\,dt,
 \qquad \mathcal B(f_0,f_1,f_2,f_3)=\int f_0T(f_1,f_2,f_3).
\]
Fix $j$ and choose $(b,c)$ by
\[
 j=1:\ (b,c)=(3,2);\qquad j=2:\ (b,c)=(3,1);\qquad
 j=3:\ (b,c)=(1,2).
\]
Only terminal labels 1 and 3 are needed; no relabeling of unequal polynomial degrees is used.

\begin{lemma}[A full neighborhood of restricted exponents]\label{new:restricted-neighborhood}
For indicator inputs of finite-measure sets in $\R^3$, $\mathcal B$ obeys a restricted
bound at the reciprocal-exponent vector $v$, where
\[
 v_0=\tfrac16,\qquad v_j=\tfrac12,\qquad v_b=\tfrac{11}{60},\qquad v_c=\tfrac16,
\]
with a finite absolute constant. It also obeys all restricted bounds at the corners
$\beta+\epsilon\sigma$, $\sigma\in\{-1,1\}^4$, with a common finite constant, where
\[
 \beta_0=\tfrac16,\qquad\beta_j=\tfrac12,\qquad
 \beta_b=\tfrac7{40},\qquad\beta_c=\tfrac16,\qquad\epsilon=\tfrac1{10000}.
\]
\end{lemma}
\begin{proof}
Lift a set in $\R^3$ to $\R^6$ using the last coordinate in each of the three blocks.
For $E_0$ use $[-1,1]$ in each of the three unused coordinates; for each of the other
sets use $[-2,2]$ in those coordinates. For $t\in[0,1]$, a point in the inner unused-coordinate
box remains in the outer one after every shift. The lifted incidence is therefore exactly
$8$ times the three-dimensional incidence. All four lifted set measures are fixed positive
multiples of the original measures. Apply \cref{new:incidence-consequence}, absorbing these
fixed multiples into the constant. This proves the bound at $v$.

The trivial bound is $\mathcal B(\one_{E_0},\ldots,\one_{E_3})\leq|E_i|$ for each $i$.
Let $w=(1/6,1/2,1/6,1/6)$ in the order $(0,j,b,c)$. Then $\sum w_i=1$ and
$\beta=(v+w)/2$. For $\beta'=\beta+\epsilon\sigma$ define
\[
 \lambda=60\left(\sum_i\beta_i'-1\right),\qquad r_i=\beta_i'-\lambda v_i.
\]
Since $\sum_i\beta_i=121/120$, $\sum_i v_i=61/60$, and $|\sum_i\sigma_i|\leq4$,
\[
 |\lambda-\tfrac12|\leq240\epsilon,
 \qquad r_i\geq\tfrac1{12}-121\epsilon>0,
 \qquad \sum_i r_i=1-\lambda.
\]
Thus $0<\lambda<1$ and $\beta'=\lambda v+\sum_i r_i e_i$ is a convex combination.
Raise the restricted bound at $v$ to the power $\lambda$ and each trivial bound at $e_i$
to $r_i$, and multiply. The powers on the left sum to one; the powers of each set measure
on the right sum to $\beta'_i$. Increasing the original constant to at least one makes the
constant uniform in all 16 corners.
\end{proof}

\begin{lemma}[Restricted-to-strong summation with no unstated interpolation theorem]
\label{new:restricted-strong}
The form $\mathcal B$ satisfies
\[
 |\mathcal B(f_0,f_1,f_2,f_3)|\leq C\|f_0\|_6\|f_j\|_2
                                      \|f_b\|_{40/7}\|f_c\|_6.
\]
\end{lemma}
\begin{proof}
Positivity permits replacing all inputs by their absolute values. By homogeneity assume
$\|f_i\|_{p_i}=1$, where $p_i=1/\beta_i$; zero norms give zero by Tonelli.
Use disjoint sets $E_{i,n}=\{2^n\leq f_i<2^{n+1}\}$, and let $m_{i,n}=|E_{i,n}|$.
They have finite measure, and
\[
 \sum_n2^{np_i}m_{i,n}\leq1,\qquad
 f_i\leq\sum_n2^{n+1}\one_{E_{i,n}}.
\]
For each tuple of nonzero measures, select from \cref{new:restricted-neighborhood} the
corner whose sign in coordinate $i$ minimizes $m_{i,n_i}^{\epsilon\sigma_i}$.
Its restricted bound is
\[
 C\prod_i m_{i,n_i}^{\beta_i}
                   \min\{m_{i,n_i}^{\epsilon},m_{i,n_i}^{-\epsilon}\}.
\]
The positive expansion of the form and Tonelli therefore bound the normalized form by
\[
 16C\prod_{i=0}^3\left(\sum_{n:m_{i,n}>0}
       2^n m_{i,n}^{\beta_i}
              \min\{m_{i,n}^{\epsilon},m_{i,n}^{-\epsilon}\}\right).
\]
For $n\geq0$, $m_{i,n}\leq2^{-np_i}\leq1$, so its summand is at most
$2^{-np_i\epsilon}$. For $n<0$ and $m_{i,n}\geq1$, it is at most $2^{np_i\epsilon}$.
For $n<0$ and $m_{i,n}<1$, it is at most $2^n\leq2^{np_i\epsilon}$, because
$0<p_i\epsilon<1$. Each parenthesis is thus at most
\[
 \sum_{n\in\mathbb Z}2^{-|n|p_i\epsilon}
 =\frac{1+2^{-p_i\epsilon}}{1-2^{-p_i\epsilon}}<\infty.
\]
This proves the displayed estimate with an explicit finite constant. Finite partial sums
and then monotone convergence justify the argument for arbitrary nonnegative inputs;
complex inputs follow by domination. No cone-summation or Marcinkiewicz assertion is left
unproved.
\end{proof}

\begin{theorem}[Improving estimates at general scale]\label{new:improving}
With $j,b,c$ chosen as above and $N>0$,
\[
 \|A_N(f_1,f_2,f_3)\|_{6/5}
 \leq C N^{-1/20}\|f_j\|_2\|f_b\|_{40/7}\|f_c\|_6.
\]
Define the adjoint, with all conjugations specified, by
\[
 A_N^{*j}(g_0,g_b,g_c)(y)
 =\frac1N\int_0^N g_0(y-t^je_j)
        \overline{g_b(y-t^je_j+t^be_b)}
        \overline{g_c(y-t^je_j+t^ce_c)}\,dt.
\]
Then
\[
 \|A_N^{*j}(g_0,g_b,g_c)\|_2
 \leq C N^{-1/20}\|g_0\|_6\|g_b\|_{40/7}\|g_c\|_6.
\]
\end{theorem}
\begin{proof}
The unit-scale form estimate gives the unit-scale $L^{6/5}$ estimate by
\cref{new:duality}. Reflection $x\mapsto-x$ changes all three negative shifts to positive
shifts and preserves all norms. For general $N$, use
$D_Nx=(Nx_1,N^2x_2,N^3x_3)$ and substitute $t=Nu$. Its determinant is $N^6$, and
\[
 A_N(f)(D_Nx)=A_1(f_1\circ D_N,f_2\circ D_N,f_3\circ D_N)(x),
 \qquad \|f_i\circ D_N\|_{p_i}=N^{-6/p_i}\|f_i\|_{p_i}.
\]
The norm factor is
\[
 N^{6(5/6-1/2-7/40-1/6)}=N^{-1/20}.
\]
For bounded supported inputs and a bounded supported test $h$, Fubini and the translation
$y=x+t^je_j$ give the exact identity
\[
 \langle A_N(f_1,f_2,f_3),g_0\rangle
 =\langle f_j,A_N^{*j}(g_0,f_b,f_c)\rangle.
\]
Every integral is absolutely convergent, since the inputs are bounded and one spatial
support is bounded uniformly in $t\in[0,N]$. Apply the first estimate and H\"older to its
left side, and then \cref{new:duality} with output exponent two, to obtain the adjoint
estimate. Truncate absolute values and spatial supports, use positivity for the absolute
adjoint, and apply monotone convergence and domination to extend to all inputs of the
stated finite norms. This also shows that the concrete parameter integral represents the
bounded extension almost everywhere.
\end{proof}


\section{The real inverse theorem, developed in dependency order}\label{sec:internal-kosz-adjoint}
The source organization is \cite[\S\S4.2--4.5]{KoszEtAl}.  The public every-input uniformity
lemma is \emph{not} a premise of this section.  The only polynomial uniformity estimate used
inside the section is the highest-active-degree estimate proved below, including its
parameter-shifted version.  We work with all increasing degree lists contained in $\{1,2,3\}$,
not just the list $(1,2,3)$: shorter lists and lower coefficients are needed by the slicing argument.

\subsection{Fej\'er measures and the existing cube quantity}
\begin{definition}[Fej\'er measures]\label{def:local-uniformity}
For $H>0$ let $\kappa_H(h)=H^{-1}(1-|h|/H)_+$. It is a probability density on $\R$.
All independent shift variables below are integrated against the explicitly specified
uniform or Fej\'er probability measure.
\end{definition}
For calculations let
\[
 C_s(f;x,\mathbf h)=\prod_{\omega\in\{0,1\}^s}
       \mathcal C^{|\omega|}f\bigl(x+(\omega\cdot\mathbf h)e_j\bigr),
 \qquad C_0(f;x)=f(x).
\]
Then
\refstepcounter{equation}
\[
 \label{patch:cube-succ}
 C_{s+1}(f;x,\mathbf h,h)=C_s(f;x,\mathbf h)
                                  \overline{C_s(f;x+he_j,\mathbf h)}.
 \tag{\theequation}
\]
In ambient dimension $n$, with an explicit positive spatial normalizer $V$, write
\[
 Q_{s,H,V}^{j}(f)=V^{-1}\operatorname{Re}\int\!\int C_s(f;x,\mathbf h)
                                   \prod_i\kappa_H(h_i)\,d\mathbf h\,dx.
\]
The notation with unequal radii, $Q_{\mathbf L,V}^{j}$, means the same expression with
$\prod_i\kappa_{L_i}(h_i)$.  These are proof-local normalizations.  For the three-dimensional monomial system the normalization is $N^6$; when a box volume
$V$ is used to obtain bounds by one, the factor $V/N^6$ is retained explicitly.
Define $T_{H,j}u(x)=H^{-1}\int_0^H u(x+te_j)\,dt$.

\begin{lemma}[Squares, positivity, and appending a coordinate]\label{patch:fejer-squares}
The difference of two independent uniform variables in $[0,H]$ has density $\kappa_H$.
Consequently
\[
 \int\!\int u(x)\overline{u(x+he_j)}\kappa_H(h)\,dh\,dx
       =\|T_{H,j}u\|_2^2.
\]
For $s\geq1$, every order-$s$ cube power above is real and nonnegative.  If $|f|\leq1$ and $f$ is supported in a
rectangular box of volume $V$, then $0\leq Q_{s,H,V}^{j}(f)\leq1$.
If the box has $j$-th side length $S$, then
\refstepcounter{equation}
\[
 \label{patch:raise-order}
 \bigl(Q_{s,H,V}^{j}(f)\bigr)^2
       \leq(1+H/S)Q_{s+1,H,V}^{j}(f).
 \tag{\theequation}
\]
\end{lemma}
\begin{proof}
Expand the square of $T_{H,j}u$, use translation in $x$, and apply Fubini.  The density of
$a-b$ for $a,b\in[0,H]$ is $H^{-2}|[0,H]\cap([0,H]-h)|=\kappa_H(h)$.
Apply this identity to $u=C_{s-1}(f;\cdot,\mathbf h)$ for the last difference.  It proves
reality and positivity before any comparison of kernels is made.  The zero vertex gives
$|C_s(f;x,\mathbf h)|\leq\mathbf1_B(x)$, so the upper bound is one.
For \eqref{patch:raise-order}, $u=C_s$ is supported in $B$, its full integral is unchanged
by $T_{H,j}$, and $T_{H,j}u$ is supported in a box of volume $V(1+H/S)$.
Thus $|\int u|^2\leq V(1+H/S)\|T_{H,j}u\|_2^2$.  Average over the first $s$ increments,
apply Jensen, and use \eqref{patch:cube-succ}.  The same argument applies to translated
support boxes.  Translation and conjugation of $f$ preserve every integrated cube power.
\end{proof}

\begin{lemma}[Fej\'er van der Corput]\label{lem:fejer-vdc}
Let $I$ have length $N>0$, $0<H\leq N/4$, and $F$ be measurable and $1$-bounded.  Then
\[
 \left|\frac1N\int_I F(t)\,dt\right|^2
 \leq4\int\kappa_H(h)
       \left|\frac1N\int_{I\cap(I-h)}F(t+h)\overline{F(t)}\,dt\right|\,dh
       +\frac{8H}{N}.
\]
\end{lemma}
\begin{proof}
Put $W=\mathbf1_IF$ and $S(u)=H^{-1}\int_0^H W(u+a)\,da$.
Then $\int S=\int_I F$ and $S$ is supported in an interval of length $N+H$.
Cauchy--Schwarz followed by square expansion gives the stronger upper bound
$(1+H/N)\int\kappa_H(h)|N^{-1}\int_{I\cap(I-h)}F(t+h)\overline{F(t)}\,dt|\,dh$.
Use $1+H/N\leq2$.  The displayed version is therefore valid; it is the existing
\pid{fejer_vdc'} interface.  The plus sign before $8H/N$ is explicit.
\end{proof}
No unrestricted polynomial-to-uniformity assertion is included in this scalar lemma.

\begin{lemma}[Signed integrated removal]\label{patch:signed-vdc}
Let $I$ have length $N>0$ and $0<H\leq N/4$.  Suppose $V>0$, $\int|g|^2\leq V$, and $\int|F(x,t)|^2\,dx\leq V$ for the relevant
parameter values.  Assume joint measurability and the integrability supplied, in applications,
by bounded factors and their finite-support envelopes.  Set
$A=V^{-1}\int g(x)\mathbb E_{t\in I}F(x,t)\,dx$.  Then
\refstepcounter{equation}
\[
 \label{patch:signed-vdc-eq}
 |A|^2\leq2\operatorname{Re}\int\kappa_H(h)
      \frac1V\int\mathbb E_{t\in I}F(x,t)\overline{F(x,t+h)}\,dx\,dh
      +2H/N.
 \tag{\theequation}
\]
\end{lemma}
\begin{proof}
First remove exactly $g$, which is independent of $t$, by Cauchy--Schwarz in $x$:
$|A|^2\leq V^{-1}\int|\mathbb E_I F(x,t)|^2dx$.
Apply the window construction in the preceding proof with $x$ retained and integrate in $x$.
It gives $(1+H/N)$ times the real part of the averaged cut-off correlation.  This whole
averaged correlation is nonnegative because it came from an integrated square.
Its replacement of $I\cap(I-h)$ by $I$ costs at most $|h|/N$: indeed
$\int|F(x,t)\overline{F(x,t+h)}|dx\leq V$ by Cauchy--Schwarz.
Bound $1+H/N$ by two and average $|h|\leq H$.
This proves \eqref{patch:signed-vdc-eq}.  In particular, it is not obtained by deleting an
absolute value from \pid{fejer_vdc'}.
\end{proof}

\subsection{The elementary coordinate-removal induction}
\begin{lemma}[Cylinder Cauchy--Schwarz, with its state]\label{patch:cylinder}
Let $z$ be an auxiliary probability-space variable and $b_i$ independent probability-space
variables, $1\leq i\leq r$.  Let $v(z,\mathbf b)$ and $u_i(z,\mathbf b)$ be jointly measurable and $1$-bounded,
with $u_i$ independent of $b_i$.  Then
\[
 \left|\mathbb E_{z,\mathbf b}v(z,\mathbf b)\prod_{i=1}^r u_i(z,\mathbf b)\right|^{2^r}
 \leq\mathbb E_{z,\mathbf b^0,\mathbf b^1}
       \prod_{\omega\in\{0,1\}^r}\mathcal C^{|\omega|}v(z,\mathbf b^\omega).
\]
The right-hand side is real and nonnegative.  The same statement for finite nonprobability
spaces follows by writing every normalization factor explicitly.
\end{lemma}
\begin{proof}
After $a$ steps use independent pairs $b_i^0,b_i^1$ for $i\leq a$ and one variable $b_i$
for $i>a$.  Define
\[
 v_a=\prod_{\omega\in\{0,1\}^a}\mathcal C^{|\omega|}v(z,\mathbf b^\omega),
 \qquad
 u_{i,a}=\prod_{\omega\in\{0,1\}^a}\mathcal C^{|\omega|}u_i(z,\mathbf b^\omega)
 \quad(i>a),
\]
and $J_a=\mathbb E[v_a\prod_{i>a}u_{i,a}]$.
At step $a+1$ the removed block is exactly $u_{a+1,a}$.
It is independent of the only variable being doubled, $b_{a+1}$.
Write the average as an outer average of this block times the inner $b_{a+1}$ average.
Cauchy--Schwarz in all the outer variables and expansion of the inner square yield
$|J_a|^2\leq J_{a+1}$.  The new $J_{a+1}$ is nonnegative because it is that square average.
The finite-product equivalence $\{0,1\}^{a+1}\simeq\{0,1\}^a\times\{0,1\}$ proves
that the expanded factors are precisely the displayed next state; the new bit toggles
conjugation.  Induct on $a=0,\ldots,r-1$.  The auxiliary variable $z$ is never doubled.
\end{proof}
This is the reusable, fully specified Gowers--Cauchy--Schwarz induction used below.
It does not claim that a polynomial family is already a cube.

\subsection{Scalar polynomial oscillation and coefficient separation}
\begin{lemma}[Real polynomial oscillation]\label{lem:real-polynomial-oscillation}
For $1\leq d\leq3$ there is $C_d\geq1$ such that
\[
 \left|N^{-1}\int_0^N\ee{\sum_{r=1}^d a_rt^r}\,dt\right|
 \leq C_d\min\{1,(\max_r|a_r|N^r)^{-1/d}\}.
\]
When $\max_r|a_r|N^r=0$, interpret the minimum as one; in Lean use a separate zero case,
not a negative real power of zero.  Hence a lower bound $0<\varepsilon\leq1$ for this average bounds every nonconstant coefficient
by $|a_r|N^r\leq (C_d/\varepsilon)^d$.
\end{lemma}
\begin{proof}
Scale to $N=1$.  For a real polynomial $p$ of degree at most $d-1$ with a coefficient of
modulus at least $A$, a measurable sublevel set $E\subseteq[0,1]$ on which $|p|\leq\theta$
has length at most $C_d(\theta/A)^{1/(d-1)}$ when $d\geq2$.
To see this, choose $d$ points in $E$ with mutual separations at least $|E|/(2d)$ by successively
removing intervals about the previously chosen points (or use $d$ separated quantile blocks).
Lagrange interpolation then bounds every coefficient by $C_d\theta|E|^{-(d-1)}$.
Null sets and degree-zero polynomials are handled separately.
Apply this to $p=P'$ and $A=\max|a_r|$.  Split $[0,1]$ at the finitely many roots of
$P''$, $P'-\theta$, and $P'+\theta$.  On each remaining interval where $|P'|\geq\theta$,
$P'$ is monotone; integration by parts in $\ee{P}/(2\pi iP')$ bounds the integral by
$C_d/\theta$.  The total number of intervals is bounded in terms of $d$.
Choose $\theta=A^{1/d}$ when $A\geq1$.  The sublevel contribution and the integration-by-parts
contribution are both $O_d(A^{-1/d})$.  For $d=1$ integrate the exponential directly;
for $A<1$ use the bound one.  Scale back.
\end{proof}

\begin{lemma}[Distinct degrees give a triangular frequency bound]\label{patch:triangular}
For admissible $P_1,\ldots,P_k$, if
$|\mathbb E_{t\in[0,N]}\ee{\sum_i\xi_iP_i(t)}|\geq\delta^{O(1)}$, then
$\sum_iN^{d_i}|\xi_i|\leq\delta^{-O(1)}$.
\end{lemma}
\begin{proof}
Write $a_r=[t^r]\sum_i\xi_iP_i(t)$ and $z_i=N^{d_i}|\xi_i|$.
The preceding lemma bounds every $N^r|a_r|$ by a power of $\delta^{-1}$.
Starting at $i=k$ and descending, use the exact identity
\[
 \xi_i\operatorname{lc}(P_i)=a_{d_i}-\sum_{u>i}\xi_u[t^{d_i}]P_u.
\]
Consequently $z_i\leq u_c(\delta)\{M+u_c(\delta)\sum_{u>i}z_u\}$ for a known power bound $M$.
There are at most three steps.  Apply \cref{patch:budgets} to this explicit recursion.
\end{proof}


\subsection{The protected highest-degree PET argument}\label{patch:pet-producer}
This subsection supplies, rather than assumes, the polynomial estimate needed inside the
real inverse proof.  It is deliberately a highest-\emph{active}-degree result.  General selected
indices will be handled later by the independent slicing argument.

\subsubsection{Symbolic state and a leading-coefficient invariant}
Use $t$ for the averaging variable and $\mathbf h=(h_1,\ldots,h_r)$ for formal indeterminates.
Polynomials are elements of $(\mathbb R[\mathbf h])^n[t]$; degrees in this subsection are degrees
in $t$, \emph{before} evaluating the $h_i$.  A normal state consists of a finite ordered list
\[
 Q_1,\ldots,Q_L,\qquad Q_i(0,\mathbf h)=0,
\]
of nonzero distinct vector polynomials, with $Q_1$ of maximum degree, together with a bounded
function block $g_i$ for each $Q_i$ and a spatial block $g_0$.
A block is a finite list of translated/conjugated input factors.  Multiplicities are retained.
The distinguished block $g_1$ is required to contain \emph{one} translated or conjugated copy
of the protected input $f$, not a product containing $f$.

Fix the protected input's original degree $D\in\{2,3\}$, leading scalar coefficient $a\ne0$,
and coordinate $e_m$.  Retain the following invariant for $Q_1$ and for every nonzero
$Q_1-Q_i$, including $Q_0=0$:
\refstepcounter{equation}
\[
 \label{patch:lead-invariant}
 \operatorname{lc}_t R=a\,p_R(\mathbf h)e_m,
 \quad p_R\in\mathbb Z[\mathbf h]\setminus\{0\},
 \quad p_R\text{ multilinear and homogeneous of degree }D-\deg_t R.
 \tag{\theequation}
\]
Only the leading coefficients have this form.  The full polynomials may contain lower
coefficients in any active coordinate and with arbitrary admissible real values.

Initially the input polynomials are $P_i(t)e_i$ with increasing degrees and protected index $m$.
Then \eqref{patch:lead-invariant} holds with $p_R=1$ for the head and its differences with
lower-degree inputs.  It also holds for the shifted family
\refstepcounter{equation}
\[
 \label{patch:initial-shift-family}
 P_i(t+\omega\cdot\mathbf h)e_i-P_i(\omega\cdot\mathbf h)e_i,
       \qquad \omega\in\{0,1\}^{r_0},\quad r_0\leq4,
 \tag{\theequation}
\]
when the protected head is the highest input with $\omega=(1,\ldots,1)$.
For two shifts of that input, their difference has degree $D-1$ and leading coefficient
$aD[(\mathbf1-\omega)\cdot\mathbf h]e_m$; differences with lower inputs still have degree $D$.
Repeated normalized polynomials among nuisance blocks are grouped.  Since $D\geq2$, the
protected head has no duplicate in this initial family.

\begin{lemma}[A labelled update preserves the protected head]\label{patch:pet-update}
Suppose the current head degree is $d\geq2$.  For a pivot $p\in\{1,\ldots,L\}$ introduce one
fresh indeterminate $u$ and form the ordered children
\[
 Q_i(t+u)-Q_p(t),\quad Q_i(t)-Q_p(t)\qquad(1\leq i\leq L),
\]
putting the shifted head first.  Subtract each child's constant term in $t$, absorb that
constant as a translation of its function, collect zero polynomials in the new spatial block,
and multiply the function lists attached to identical remaining polynomials.
If the shifted head has maximum degree, the new state is normal, its head block is a
singleton of origin $f$, and \eqref{patch:lead-invariant} is preserved.
\end{lemma}
\begin{proof}
Let $\delta_u Q_1=Q_1(t+u)-Q_1(t)$.  Its degree is $d-1$ and leading coefficient
$d u\operatorname{lc}_tQ_1$.  For every $i$, compare the new head with an unshifted child:
\[
 [Q_1(t+u)-Q_p(t)]-[Q_i(t)-Q_p(t)]
       =\delta_uQ_1+(Q_1-Q_i).
\]
If the degrees of the two summands differ, take the larger.  If they agree, the leading
coefficient is the sum of an old polynomial in $\mathbf h$ and $du$ times another nonzero
old polynomial.  It cannot vanish because $u$ is a fresh indeterminate.  The sum is still
$a$ times a nonzero integer multilinear polynomial with the asserted homogeneous degree.
For a shifted child the difference is simply $(Q_1-Q_i)(t+u)$, whose leading coefficient is
unchanged.  Apply the first calculation with $i=p$ to the new head itself.
Subtracting constants does not change these nonconstant leading terms.
In particular the head differs nontrivially from every other child, and is nonconstant even
when $p=1$.  Therefore grouping never multiplies another factor into the protected head.
The shifted child conjugates its function, whereas an unshifted child retains its conjugation
flag.  Translation by its constant term preserves origin and singleton status.
\end{proof}
If all integer polynomials in the invariant have coefficient $\ell^1$-norm at most $M$,
the proof gives the bound $(D+1)M$ after an update.  Initially one can take
$M=1+2Dr_0$.  These are numerical bounds independent of the values of the real lower coefficients.

\subsubsection{The pivot is specified, and termination has a uniform bound}
Let $w_q$ count distinct leading coefficients of degree-$q$ polynomials in the normal state.
Order $(w_D,\ldots,w_1)$ lexicographically.  Stop when every degree is at most one.
Otherwise let $d=\deg_tQ_1$ and let $l$ be the minimum degree present.  Choose the first
index in the ordered list satisfying the relevant case:
If $l<d$, take the first $p$ with $\deg Q_p=l$.
If $l=d$ and $w_d>1$, take the first $p$ with
$\operatorname{lc}Q_p\ne\operatorname{lc}Q_1$.
If $l=d$ and $w_d=1$, take $p=1$.
For the first two cases the head retains maximum degree; for the third it has degree $d-1$
and every child has degree at most $d-1$.  For the selected degree $l$, the new leading
classes are the nonzero differences with the pivot leading coefficient.  Thus
\[
 w'_q=w_q\ (q>l),\qquad w'_l=w_l-1,\qquad 0\leq w'_q\leq2L\ (q<l).
\]
The type strictly decreases and the new length is at most $2L$.
The argument is in the formal polynomial ring; numerical exceptional parameter values do
not change the choice of pivot.

For an actual uniform iteration bound define $B(w,L)$ by well-founded recursion on the type,
with its value a function of the length budget $L$:
\[
 B(w,L)=0\quad\hbox{at a linear type},\qquad
 B(w,L)=1+\max_{w'\in\mathscr C(w,L)}B(w',2L)\quad\hbox{otherwise},
\]
where $\mathscr C(w,L)$ is the finite set of types with the preceding displayed bounds for
the chosen $l$.  Each $w'<w$, so this is a well-founded definition.  The usual induction
proves that it bounds the execution on every family of length at most $L$ and type $w$.
If the initial length is at most $L_0\leq3\cdot2^4$, take the finite maximum of $B(w,L_0)$
over $0\leq w_q\leq L_0$.  Call it $R_0$.  The terminal length is at most $L_1=2^{R_0}L_0$.
This proves uniform complexity, not merely termination of each individual execution.
For definiteness, $l$ in this recursion is the least $q$ with $w_q>0$; the maximum head
degree is the greatest such $q$. In the recursive maximum take all tuples satisfying the
displayed bounds, including tuples not realized by an actual family. This set is finite:
each of the at most three entries is either fixed or lies in $\{0,\ldots,2L\}$.
Every recursive call lowers the lexicographic tuple of natural numbers, a well-founded
order on a fixed finite product. To prove the bound, induct in that well-founded order,
then apply the induction hypothesis to the actual child tuple and its length bound $2L$.
Thus it bounds every possible grouping of lower-coefficient coincidences, not just one
chosen generic family.


\paragraph{Full coefficient and translation bounds.}
For a fixed spatial coordinate $i$, retain a bound $M N^{d_i-q}$ for the coefficient of
$t^q$ of every current polynomial and constant-translation record; coefficients with
$q>d_i$ remain zero.  A shift $t\mapsto t+u$, $|u|\leq H\leq N$, transforms that
coefficient into $\sum_{v\geq q}\binom vq c_v u^{v-q}$, whose modulus is at most
$2^{d_i}M N^{d_i-q}$.  Subtracting the pivot multiplies this bound by at most two.
Adding a constant translation to a function record adds two such bounds at $q=0$.
Induction for the at most $R_0$ steps gives a fixed numerical multiple of the initial budget.
This proves every support-envelope and coefficient upper bound used by the analytic execution.

\paragraph{Lean state.}
Use an outer polynomial in $t$ with coefficients in a multivariate polynomial ring in the
shift variables.  A factor carries its origin index, conjugation flag, and constant-translation
polynomial.  Normalization groups \emph{function lists}; it does not discard duplicate factors.
The head list has length one.  Polynomial equality can be handled classically; computability
of comparisons between arbitrary real coefficients is not asserted.  The termination proof
uses the lexicographic well-founded relation, not an incorrect induction on maximum degree alone.

\subsubsection{The analytic update, with no assumed cube}
For fixed old shifts let
\[
 A=V^{-1}\int g_0(x)\mathbb E_{t\in[0,N]}
                   \prod_{i=1}^L g_i(x+Q_i(t))\,dx,
\]
where every block is $1$-bounded and has $L^2$ norm squared at most $V$.
Apply \cref{patch:signed-vdc} with $g=g_0$ and the remaining product as $F$.
Then translate $x$ by $-Q_p(t)$ in the correlation on its right side.
The two children of a factor $(g_i,Q_i)$ are exactly
\[
 (\overline{g_i},Q_i(t+u)-Q_p(t)),\qquad(g_i,Q_i(t)-Q_p(t)).
\]
Normalize constants as in \cref{patch:pet-update}.  The old spatial block was removed by the
Cauchy--Schwarz inequality; constant children form the new spatial block.  This proves
\refstepcounter{equation}
\[
 \label{patch:pet-analytic-step}
 |A(\mathbf h)|^2\leq2\operatorname{Re}\mathbb E_{u\sim\kappa_H}
                       A'(\mathbf h,u)+2H/N.
 \tag{\theequation}
\]
A product block is dominated by any one of its factors, so all required $L^2$ and $L^1$
bounds remain at most $V$.  The shifted head is a translated/conjugated copy of $f$.
Every joint integral is measurable; no numerical shift is selected during this symbolic recursion.

\subsubsection{The affine endpoint: removal order and cube invariant}
At a formal linear state, \eqref{patch:lead-invariant} says that the head slope and its differences
with every other slope are multiples of $e_m$.  Hence all the slopes are multiples of $e_m$.
Include the spatial block as a nuisance of slope zero, list all nuisances in their inherited
order $0,\ldots,s-1$, and place the protected block last, at index $s$.
For fixed old shifts the quantity has the exact form
\[
 A_0=V^{-1}\int\mathbb E_{t\in[0,N]}\prod_{i=0}^{s}g_i(x+a_it e_m)\,dx,
 \qquad g_s=\hbox{one translate/conjugate of }f.
\]
For $r\leq i\leq s$ define
\[
 G_{i,0}(x)=g_i(x).
\]
\refstepcounter{equation}
\[
 \label{patch:affine-block}
 G_{i,r+1}(x;\mathbf u,u_r)=G_{i,r}(x;\mathbf u)
     \overline{G_{i,r}(x+(a_i-a_r)u_re_m;\mathbf u)}.
 \tag{\theequation}
\]
\[
 A_r(\mathbf u)=V^{-1}\int\mathbb E_t
           \prod_{i=r}^sG_{i,r}(x+a_it e_m;\mathbf u)\,dx.
\]
The following induction gives the precise cube identity. At $r=0$ the Boolean cube has
one element and the product is $g_i$. In the successor product split the vertices according
to their last bit. The last-bit-zero product is $G_{i,r}(x)$ by the induction hypothesis.
The last-bit-one product is its conjugate evaluated at $x+(a_i-a_r)u_re_m$, because the
extra bit changes the parity once. Their product is the defining formula for $G_{i,r+1}$.
Thus, for every $r$, the invariant is
\refstepcounter{equation}
\[
 \label{patch:affine-cube-invariant}
 G_{i,r}(x;\mathbf u)=\prod_{\omega\in\{0,1\}^r}\mathcal C^{|\omega|}
        g_i\left(x+\sum_{\nu<r}\omega_\nu(a_i-a_\nu)u_\nu e_m\right).
 \tag{\theequation}
\]
At step $r$ translate by $-a_rt e_m$ and remove \emph{exactly} $G_{r,r}$ by
\cref{patch:signed-vdc}.  It is independent of $t$ after that translation.
Equation \eqref{patch:affine-block} identifies the resulting correlation with $A_{r+1}$, so
\[
 |A_r|^2\leq2\operatorname{Re}\mathbb E_{u_r\sim\kappa_H}A_{r+1}+2H/N.
\]
The zero vertex bounds $|G_{i,r}|$ by $|g_i|$, and $|A_r|\leq1$.
At $r=s$ translate away $a_st e_m$; the parameter average disappears and the result is
\refstepcounter{equation}
\[
 \label{patch:affine-terminal}
 A_s(\mathbf u)=V^{-1}\int
     \prod_{\omega\in\{0,1\}^s}\mathcal C^{|\omega|}
         g_s\left(x+\sum_{\nu<s}\omega_\nu(a_s-a_\nu)u_\nu e_m\right)dx.
 \tag{\theequation}
\]
Translation/conjugation invariance identifies its integrated power with that of $f$.
The raw cube in \eqref{patch:affine-terminal} makes sense even at a parameter where a slope
difference evaluates to zero.  We only require nonzero gaps later, for changing variables.
Thus no exceptional set is hidden in the analytic recursion.

\begin{lemma}[Explicit accumulation of the removal losses]\label{patch:cs-loss}
Suppose a chain of $T$ such steps has $|A_r|\leq1$, probability shift measures, and error
$2H/N$ at each step.  If its final signed average is a nonnegative cube power $Q\leq1$, then
\[
 |A_0|^{2^T}\leq c_T\bigl(Q+T H/N\bigr),\qquad c_T=2^{2^T-1}.
\]
The same assertion holds with all earlier shift parameters retained and averaged.
\end{lemma}
\begin{proof}
Set $M_r=\mathbb E|A_r|$ except at the last stage, where use $Q=\operatorname{Re}\mathbb E A_T$.
Jensen gives $M_r^2\leq2M_{r+1}+2H/N$, with $Q$ in the last line.
For $\phi(x)=x^2/2$, $x,y\in[0,1]$ imply
$|\phi(x)-\phi(y)|\leq|x-y|$.  Inductively
$\phi^{\circ r}(M_0)\leq M_r+rH/N$, and at the last stage use $Q$.
For a completely scalar induction, split into $\phi^{\circ r}(M_0)\leq M_r$ or its reverse
and use $(x^2-y^2)/2=(x-y)(x+y)/2$ on $[0,1]$.
Since $\phi^{\circ T}(x)=x^{2^T}/c_T$, the result follows.
The last signed average is essential; it is not replaced by absolute cube correlations.
\end{proof}

\subsubsection{Real sublevel exclusion and the common-scale cube}
At the terminal stage every head gap, including the gap with the spatial slope zero, has form
\[
 a\,p_i(\mathbf h),\qquad p_i\in\mathbb Z[\mathbf h]\setminus\{0\},
\]
where $p_i$ is multilinear homogeneous of degree $D-1\leq2$, with a coefficient bound fixed
by the recursion above.  All old parameter shifts have the same radius $H$.

\begin{lemma}[The sublevel estimate actually needed here]\label{patch:multiaffine-sublevel}
Let $p\ne0$ be a multilinear integer polynomial of degree at most two.  If the coordinates
of $u$ are independently distributed with density $\kappa_1$, then, for $0<\varepsilon<1$,
\[
 \Pr\{|p(u)|\leq\varepsilon\}\leq4\sqrt\varepsilon.
\]
\end{lemma}
\begin{proof}
A nonzero constant has modulus at least one.  For a nonconstant linear polynomial, fix all
but a coordinate with a nonzero integer coefficient.  The sublevel interval has length at
most $2\varepsilon$, and $\kappa_1\leq1$.
If a quadratic term is present, write $p=A(u')u_r+B(u')$ where $A$ is linear and has a
nonzero integer coefficient.  The probability $|A|\leq\sqrt\varepsilon$ is at most
$2\sqrt\varepsilon$ by the linear case.  On its complement, the conditional sublevel
probability in $u_r$ is at most $2\varepsilon/|A|\leq2\sqrt\varepsilon$.
Integrate the remaining probability variables.  No dimension-dependent inverse theorem is used.
\end{proof}

Let $T_0$ bound the total number of parameter and affine removals; for example use
$T_0=4+R_0+L_1$.  If the original normalized lower bound is $\lambda\leq1$, choose
\[
 H/N=\min\{1/32,\lambda^{2^{T_0}}/(4T_0c_{T_0})\}>0.
\]
Take $\lambda$ to be the fixed lower bound returned from the input budget, rather than
the measured value of the particular correlation.  This choice of $H/N$ has both a power
lower bound and a power upper bound by the budget rules.  The loss lemma gives an averaged raw
terminal cube at least $\kappa=\lambda^{2^{T_0}}/(2c_{T_0})$ after harmless weakening.
For $s$ terminal gaps, take $\varepsilon=(\kappa/(16s))^2$.
Homogeneity gives $p_i(Hu)=H^{D-1}p_i(u)$.
By \cref{patch:multiaffine-sublevel}, the union of all bad gap sets has probability at most
$\kappa/4$.  Let $F$ be the raw cube after averaging all the new affine increments, as a function
of the old shifts. Then $0\leq F\leq1$ by the last-square identity, also when a gap is zero.
Its mean is at least $\kappa$. If every old shift outside the bad union had
$F<\kappa/2$, then splitting the probability integral would give
\[
 \mathbb E F\leq\Pr(\text{bad union})+\kappa/2\leq3\kappa/4<\kappa,
\]
a contradiction. Hence there are old shifts outside the union with $F\geq\kappa/2$.
Only this existence assertion is used; a measurable choice of old shifts is unnecessary.  At those shifts
\[
 |a|H^{D-1}\varepsilon\leq |a p_i(\mathbf h)|
       \leq |a|H^{D-1}M_1,
\]
where $M_1$ is the integer coefficient $\ell^1$ bound.
Changing the new affine increments by $z_i=a p_i(\mathbf h)u_i$ therefore produces
Fej\'er radii $L_i=|a p_i(\mathbf h)|H$ with
\refstepcounter{equation}
\[
 \label{patch:pet-physical-radii}
 \delta^{O(1)}N^D\leq L_i\leq\delta^{-O(1)}N^D.
 \tag{\theequation}
\]
Here $D$ is the protected polynomial degree, not the sum of ambient degrees.

\begin{lemma}[Mixed radii to one radius, with positivity before comparison]\label{patch:uniformize}
Suppose $f$ is $1$-bounded on a box of volume $V$ and $j$-th side length $S$.
If $L_i>0$ and $L\geq2\max_iL_i$, then
\[
 (Q_{\mathbf L,V}^{j}(f))^2
 \leq\left(1+\frac LS\right)\prod_i\frac{2L}{L_i}\ Q_{s+1,L,V}^{j}(f).
\]
\end{lemma}
\begin{proof}
For fixed $z_1,\ldots,z_s$ put $u_z=C_s(f;\cdot,z)$.
The zero vertex supports $u_z$ on the original box, so
\[
 |V^{-1}\!\int u_z|^2\leq(1+L/S)\,P(z),\qquad
 P(z)=V^{-1}\|T_{L,j}u_z\|_2^2\geq0.
\]
Since $L_i\leq L/2$, pointwise
$\kappa_{L_i}(z_i)\leq(2L/L_i)\kappa_L(z_i)$.
Apply this comparison only to the nonnegative $P$, after Jensen in the first display.
Square expansion of $P$ appends one cube coordinate by \eqref{patch:cube-succ}; all $s+1$
increments now have the same Fej\'er density.  This proves the identity and inequality claimed.
\end{proof}
Choose a deterministic $L=u_C(\delta)N^D$ exceeding twice the upper bounds in
\eqref{patch:pet-physical-radii} and comparable, up to a power of $\delta$, to the protected
support side length.  Every factor in \cref{patch:uniformize} is power-controlled.
If needed, pad the resulting order to a single fixed order using \eqref{patch:raise-order}.
This proves highest-active-degree uniformity control with polynomial losses, including the
shifted families \eqref{patch:initial-shift-family}; its protected factor is still the original $f$.

\subsubsection{Polynomial phases: the four initial steps are also explicit}
Suppose the correlation also contains $\ee{p_x(t)}$, with $p_x$ a measurable-in-$x$ polynomial
in $t$ of degree at most three.  Its coefficients need not be bounded.
Before any polynomial-dependent translation of $x$, apply \cref{patch:signed-vdc} four times,
retaining $x$ as a fixed outer variable.  After $r$ steps the parameter-dependent part is
\[
 \prod_{i=1}^{m}\prod_{\omega\in\{0,1\}^r}\mathcal C^{|\omega|}
     f_i(x-P_i(t+\omega\cdot\mathbf h)e_i)\;
 \ee{\sum_{\omega\in\{0,1\}^r}(-1)^{|\omega|}p_x(t+\omega\cdot\mathbf h)}.
\]
The initial step removes the original spatial factor.  At later steps use the indicator of
a fixed enlarged spatial box: the zero-shift vertex supports the integrand there for
$t\in[0,N]$.  Polynomial coefficient bounds and $4H\leq N/8$ give a single power-controlled
volume for this box.  The indicator is removed by the same Cauchy--Schwarz step.
At $r=4$ the phase is identically one, by the fourth finite-difference identity for cubics.
Normalize the constant terms of the remaining polynomial shifts and run the preceding PET
argument, retaining all four old shift variables in its averages and sublevel exclusions.

There is one degenerate initial case to treat separately.  If the highest active degree is
one, there is only one active input.  After the four steps, translating $x$ by its affine
$P_1(t)$ removes $t$ and gives the raw order-four cube of $f_1$ with increments
$\operatorname{lc}(P_1)h_i e_1$ directly.  Keep the signed average at this last step.
Its radius is $|\operatorname{lc}(P_1)|H>0$ and is comparable to $N$ up to powers of $\delta$.
Uniformize and pad as above.  Do \emph{not} group these equal-slope head copies and call the
result a singleton head.

\begin{proposition}[Highest-active-input control, including measurable polynomial phases]
\label{patch:highest-control}
Fix $1\leq d_1<\cdots<d_k\leq3$, $1\leq m\leq k$, and an integer budget $c\geq1$.
There are integers $s_0\geq2$ and $C\geq c$, depending only on these data, with the following
property. Let $0<\delta\leq1$, $N>0$, let $P_i$ for $i\leq m$ be admissible with budget
$c$, and let $f_0,f_1,\ldots,f_m$ be Borel, bounded by one, and supported in the budget-$c$
box in $\R^k$. If $p_x(t)$ is a polynomial in $t$ of degree at most three, with arbitrary
Borel real coefficients in $x$, and
\[
 \left|\int f_0(x)\mathbb E_{t\in[0,N]}\prod_{i=1}^m
          f_i(x-P_i(t)e_i)\ee{p_x(t)}\,dx\right|\geq\ell_c(\delta)N^{D_k},
\]
then, at the deterministic radius $H=u_C(\delta)N^{d_m}$,
\[
 Q_{s_0,H,N^{D_k}}^m(f_m)\geq\ell_C(\delta).
\]
The order can be chosen uniformly over the input functions, lower coefficients, and phase
coefficients. Coordinates $m+1,\ldots,k$ are passive spatial variables throughout the proof.
\end{proposition}
\begin{proof}
The four phase steps, the normal-state update with invariant \eqref{patch:lead-invariant},
the explicit pivot recursion, \eqref{patch:affine-terminal}, the loss lemma, the sublevel
exclusion, and \cref{patch:uniformize} give the result.  Each ingredient has been proved above.
The degree-one branch is the direct cube calculation just given.  Dividing by the chosen
spatial volume and finally converting to the stated box normalization costs only a budgeted
factor; the exact ambient factor remains $N^{\sum d_i}$.
\end{proof}
The proposition is not an extra hypothesis in any later theorem.


\subsection{Fourier selection and the two difference lemmas}\label{patch:auxiliary-differences}
These are the real versions of the preparation tools in \cite[\S4.5.1]{KoszEtAl}.
Their Cauchy--Schwarz steps are expanded here using \cref{patch:cylinder}.
For a scalar function supported on an interval $J$ of length $S>0$, define the paired cube
\[
 D_{a,b}^{r}f(x)=\prod_{\omega\in\{0,1\}^r}\mathcal C^{|\omega|}
       f\left(x+\sum_{i=1}^r[\omega_i a_i+(1-\omega_i)b_i]\right),
 \quad a,b\in[0,H]^r.
\]
After translation in $x$, averaging this expression over $a,b$ is the same as averaging
$C_r$ against $r$ independent Fej\'er densities.  All scalar normalizers in this subsection
are written as $S$; changing to a larger interval costs its explicit length ratio.

\begin{lemma}[One-dimensional $U^2$ inverse estimate and measurable frequencies]
\label{patch:u2-fourier-selection}
If $|f|\leq1$ and $f$ is supported in $J$, then
\[
 Q_{2,H,S}(f)\leq H^{-2}\|\widehat f\|_\infty^2.
\]
For a jointly measurable bounded family with uniformly bounded support, a power lower bound
for these powers on a measurable set of parameters gives a measurable choice of frequencies
with Fourier coefficients of power size.  A finite frequency set is not assumed.
\end{lemma}
\begin{proof}
Put $g_a(x)=f(x)\overline{f(x+a)}$.  By \cref{patch:fejer-squares},
$SQ_{2,H,S}=\int\kappa_H(a)\|T_Hg_a\|_2^2da$.
Use $\kappa_H\leq H^{-1}$ against this nonnegative integrand and apply \cref{new:integral-plancherel} in $x$.
For each $\xi$, a second application of \cref{new:integral-plancherel}, this time in $a$, gives
\[
 \int|\widehat{g_a}(\xi)|^2da
    =\int|\widehat f(\zeta+\xi)|^2|\widehat f(\zeta)|^2d\zeta
    \leq\|\widehat f\|_\infty^2\|f\|_2^2.
\]
Here the second identity can be verified without an unproved ambiguity-function formula.
For fixed $\xi$ put $b_\xi(a)=\int f(x)\overline{f(x+a)}\ee{-\xi x}\,dx$.
Bounded support makes $b_\xi$ integrable, square-integrable, and compactly supported in the
$a$ variable. Fubini gives
\[
 \widehat {b_\xi}(\zeta)=\widehat f(\xi-\zeta)\,
                                  \overline{\widehat f(-\zeta)}.
\]
Indeed, set $y=x+a$ in its defining double integral; its absolute integral is at most
$\|f\|_1^2$. Applying the integral Plancherel identity to $b_\xi$ and then changing
$\zeta$ to $-\zeta$ proves exactly the preceding display.
The multiplier of $T_H$ at $\xi$ is $\widehat w_H(-\xi)$, where
$w_H=H^{-1}\mathbf1_{[0,H]}$. Its squared $L^2$ norm is $H^{-1}$ by the integral
Plancherel identity and reflection invariance.  Hence $SQ_2\leq H^{-2}\|\widehat f\|_\infty^2\|f\|_2^2$;
use $\|f\|_2^2\leq S$.
For selection enumerate $\mathbb Q$.  The Fourier transform of each integrable section is
continuous in frequency, and its value at each rational is measurable in the parameter by
Fubini.  Where the supremum exceeds a threshold $a$, choose the first rational at which the
modulus exceeds $a/2$.  The disjoint least-index sets are measurable and cover the required
parameter set.  Outside it use zero.  Popularity supplies a uniform positive threshold
before selection.  This proves the claimed jointly measurable choice and does not infer a
uniform finite set from continuity alone.
\end{proof}

\begin{lemma}[Dual--difference interchange, explicit cylinder proof]\label{patch:dual-difference}
Let $F(x)=\mathbb E_tF_t(x)$, where the $F_t$ are jointly measurable, $1$-bounded and supported
in a common interval $J$.  Let $E\subseteq[0,H]^{2r}$ and $\varphi(a,b)$ be measurable, $r\geq1$.
Put
\[
 \theta=\mathbb E_{a,b\in[0,H]^r}\mathbf1_E(a,b)
           \left|S^{-1}\int D_{a,b}^{r}F(x)\ee{\varphi(a,b)x}\,dx\right|^2.
\]
Write $J_r=J-[0,rH]$, $T=|J_r|/S$, and
\[
 \Box^r E=\{(a,b,c): (a,b^\omega)\in E\text{ for every }\omega\in\{0,1\}^r\},
 \quad b_i^\omega=(1-\omega_i)b_i+\omega_i c_i.
\]
Then, with $G_{b,c}(x)=\mathbb E_tD_{c,b}^{r}F_t(x)$ and
$\psi(a,b,c)=\sum_\omega(-1)^{|\omega|}\varphi(a,b^\omega)$,
\refstepcounter{equation}
\[
 \label{patch:dual-diff-bound}
 \mathbb E_{a,b,c}\mathbf1_{\Box^r E}
       \left|S^{-1}\int G_{b,c}(x)\ee{\psi(a,b,c)x}\,dx\right|^2
 \geq\theta^{2^r}T^{-2(2^r-1)}.
 \tag{\theequation}
\]
\end{lemma}
\begin{proof}
Expand the square defining $\theta$ using two spatial variables $x,z$ and expand every copy
of $F$ using independent parameters $t_{\omega,0},t_{\omega,1}$.
The vertex $\omega=0$ contributes
\[
 v=\mathbf1_E(a,b)F_{t_{0,0}}(x+\textstyle\sum b_i)
       \overline{F_{t_{0,1}}(z+\textstyle\sum b_i)}\,
       \ee{\varphi(a,b)(x-z)}.
\]
Partition the other Boolean vertices by their first nonzero coordinate:
$\Omega_i=\{\omega:\omega_1=\cdots=\omega_{i-1}=0,\ \omega_i=1\}$.
The product of their two spatial factors for $\omega\in\Omega_i$ is a $1$-bounded nuisance
$u_i$ independent of $b_i$.  Normalize $x,z$ on $J_r$, and normalize all parameter variables
by their probability measures.  The initial average is $\theta/T^2$.
Apply \cref{patch:cylinder}, doubling $b_1,\ldots,b_r$ in that order.
At stage $i$ it removes exactly the block indexed by $\Omega_i$ and hence all its unused
$t$ parameters.  The only parameters surviving at the endpoint are $t_{0,0},t_{0,1}$.
The indicator becomes the product over all corners, namely $\mathbf1_{\Box^r E}$;
the phase becomes $\ee{\psi(x-z)}$; the two function products become
$D_{c,b}^{r}F_{t_{0,0}}(x)$ and its conjugate at $z$.
Integrating $t_{0,0},t_{0,1},x,z$ is exactly the squared coefficient of $G_{b,c}$.
The support remains in $J_r$, so these spatial integrals may be written over $\mathbb R$.
Multiply the cylinder inequality by $T^2$ to obtain \eqref{patch:dual-diff-bound}.
\end{proof}

\begin{lemma}[Removal of missing-coordinate phases]\label{patch:missing-phase}
Let $f$ be as above and let $\varphi_i(a,b)$ be independent of $b_i$, $1\leq i\leq r$.
If
\[
 \theta=\mathbb E_{a,b}
       \left|S^{-1}\int D_{a,b}^{r}f(x)\ee{(\sum_i\varphi_i(a,b))x}\,dx\right|^2,
\]
then
\refstepcounter{equation}
\[
 \label{patch:missing-phase-bound}
 Q_{r+1,H,S}(f)\geq
       \frac{\theta^{2^r}}{T^{2(2^r-1)}(1+H/S)},\qquad T=(S+rH)/S.
 \tag{\theequation}
\]
\end{lemma}
\begin{proof}
Expand the square using $x,z\in J_r$.
Use as head $v=f(x+\sum b_i)\overline{f(z+\sum b_i)}$.
For nuisance $u_i$, multiply the vertex factors indexed by $\Omega_i$ from the preceding
proof by $\ee{\varphi_i(a,b)(x-z)}$.  This block is independent of $b_i$.
The same cylinder induction gives
\[
 R:=\mathbb E_{b,c}\left|S^{-1}\int D_{c,b}^{r}f(x)\,dx\right|^2
       \geq\theta^{2^r}T^{-2(2^r-1)}.
\]
For each $b,c$, $u=D_{c,b}^{r}f$ is supported in a translate of $J$, by its zero vertex.
Its full integral equals $\int T_Hu$, and the latter function has support length at most $S+H$.
Thus $|S^{-1}\int u|^2\leq(1+H/S)S^{-1}\|T_Hu\|_2^2$.
Average in $b,c$ and expand this last square.  It is exactly the paired cube with $r+1$
uniform pairs in $[0,H]$, hence $Q_{r+1,H,S}$.  This proves \eqref{patch:missing-phase-bound}.
\end{proof}
These two lemmas are useful even when the support/shift ratios are large: the displayed
ratios are then retained and bounded by the budget rules.  No comparison of complex
integrands by positive kernels is implicit.

\begin{lemma}[A measurable dummy phase with a small exceptional set]\label{patch:dummy-phase}
For $r\geq1$, $H,M>0$, and $0<\epsilon<1$, there is a finite-valued measurable function
$\Phi(a,b)$ on $[0,H]^{2r}$ such that its alternating sum
$\Psi(a,b,c)=\sum_\omega(-1)^{|\omega|}\Phi(a,b^\omega)$ satisfies $|\Psi|>M$ except on
a set of relative measure at most $\epsilon$ in $[0,H]^{3r}$.
\end{lemma}
\begin{proof}
Choose $n$ with $r2^{-n}\leq\epsilon$ and partition each $[0,H]$ into $2^n$ half-open
intervals, with the final endpoint assigned to the last cell.
Enumerate the resulting $r$-dimensional cells by distinct nonnegative integers $\iota$.
On a cell of the $b$ variable put $\Phi(a,b)=2M\,2^{\iota(b)}$.
Except when some $b_i,c_i$ are in the same one-dimensional cell, the $2^r$ corners $b^\omega$
are in distinct cells.  A signed sum of distinct powers of two has absolute value at least
one: its largest term exceeds the sum of all smaller possible terms.
Thus $|\Psi|\geq2M$ there.  The union of the $r$ bad coordinate events has relative measure
at most $r2^{-n}$.  Boundaries are null.  This grid resolves the purpose-built function
$\Phi$, not the level sets of an arbitrary input.  The possibly enormous phase values are
permitted: later phase mappings are arbitrary measurable functions, and their exponentials
still have modulus one.
\end{proof}

\subsection{Conditional degree lowering with all hypotheses exposed}\label{patch:conditional-dl}
Let $\mathbb R^0$ be a point with measure one.  For fixed $k,l,m$ with $1\leq m\leq l\leq k$,
define $\mathsf{MA}(m,l)$ to be the following assertion. For every integer input budget
$c\geq1$ there is an integer $C\geq c$ with this property for every $0<\delta\leq1$
and $N>0$. The $P_i$ are admissible with budget $c$. The Borel functions
$g_0,\ldots,g_{m-1}$ on $\mathbb R^{m-1}$ are bounded by one and supported in
$\prod_{i<m}[-u_c(\delta)N^{d_i},u_c(\delta)N^{d_i}]$.
The mappings $\zeta_m,\ldots,\zeta_l$ are measurable and
$\theta_{l+1},\ldots,\theta_k$ are fixed real numbers.
A lower bound $\ell_c(\delta)N^{D_{m-1}}$ for
\refstepcounter{equation}
\[
 \label{patch:ma-correlation}
 \left|\int g_0(y)\mathbb E_t\prod_{i<m}g_i(y-P_i(t)e_i)
       \ee{\sum_{i=m}^l\zeta_i(y)P_i(t)+\sum_{i>l}\theta_iP_i(t)}\,dy\right|
 \tag{\theequation}
\]
implies that the canonical set
\refstepcounter{equation}
\[
 \label{patch:ma-set}
 Y=\left\{y\text{ in the support envelope}:\
      \sum_{i=m}^lN^{d_i}|\zeta_i(y)|+\sum_{i>l}N^{d_i}|\theta_i|
                           \leq u_C(\delta)\right\}
 \tag{\theequation}
\]
has measure at least $\ell_C(\delta)N^{D_{m-1}}$. Here the support envelope means
$\prod_{i<m}[-u_C(\delta)N^{d_i},u_C(\delta)N^{d_i}]$, allowing the output budget to
absorb the bounded enlargements used below.
There is no bound on the phase mappings before this conclusion.  In the real case no
integer denominator occurs.  The output budget is independent of the phase values.

\begin{lemma}[Conditional structured degree lowering]\label{patch:conditional-degree}
Assume the theorem $\mathsf{MA}(m,l)$ for every budget.  For $s\geq3$ and admissible data set
\[
 p_x(t)=\sum_{i=m+1}^l\xi_i(x)P_i(t)+\sum_{i>l}\eta_iP_i(t),\qquad x\in\R^m.
\]
\[
 F(x)=\mathbb E_t f_0(x+P_m(t)e_m)
           \prod_{i<m}f_i(x-P_i(t)e_i+P_m(t)e_m)
           \ee{p_{x+P_m(t)e_m}(t)}.
\]
Here the $f_i$ are $1$-bounded with budgeted support, the $\xi_i$ are measurable, and the
$\eta_i$ are fixed.  If $H$ is comparable to $N^{d_m}$ up to budgeted powers and the
order-$s$ uniformity norm of $F$ in direction $e_m$ is power-large, then its order-$(s-1)$
norm at the \emph{same} $H$ is power-large.  Enlarge the support normalizer to contain $F$;
its volume ratio is explicitly power-controlled.
\end{lemma}
\begin{proof}
Write $r=s-2$, $x=(y,z)$ with $y=x_{<m}$ and $z=x_m$.
The following steps specify every selection and every use of the assumed major-arc theorem.

\paragraph{1. Fourier coefficients on popular sections.}
Expand the first $r$ paired differences in the order-$s$ power.  Fubini identifies the
remainder with a one-dimensional $U^2$ power on each $y$ section.
Use \cref{patch:budgets,patch:u2-fourier-selection}.  There is a measurable power-large set
of $y$ and, on each such section, a measurable power-large set $E_y\subseteq[0,H]^{2r}$
on which
\[
 \left|S^{-1}\int D_{a,b}^{r}F_y(z)\ee{\varphi_y(a,b)z}\,dz\right|\geq\delta^{O(1)}.
\]
Here $S$ is the length of one fixed enlarged section interval $J$, comparable to $N^{d_m}$
up to powers.  The frequency $\varphi_y(a,b)$ is the least-rational selection described above,
so it is jointly measurable.  Translated supports of the differences are all contained in
$J-[0,rH]$ after one further budgeted enlargement if necessary.

\paragraph{2. Interchange the average and the differences.}
Write $F_y=\mathbb E_tF_{t,y}$ and apply \cref{patch:dual-difference} on each good $y$.
Define
\[
 G_{y,b,c}(z)=\mathbb E_tD_{c,b}^{r}F_{t,y}(z),\qquad
 \psi_y(a,b,c)=\sum_\omega(-1)^{|\omega|}\varphi_y(a,b^\omega).
\]
Popularity gives a canonical measurable set $Z$ of pairs $(y,k)$, $k=(a,b,c)$, on which
$(a,b^\omega)\in E_y$ for every corner and
$|\int G_{y,b,c}(z)\ee{\psi_y(k)z}dz|\geq\delta^{O(1)}N^{d_m}$.
The measure of $Z$ is power-large in the units $N^{D_{m-1}}H^{3r}$.

\paragraph{3. Interchange the section and cube parameters.}
Fubini and popularity produce a measurable set $K\subseteq[0,H]^{3r}$ of power-large relative
measure such that each $k\in K$ has a power-large section $Z_k$ in the $y$ variables.
Every later set can be defined by explicit thresholds on these joint measurable functions;
no arbitrary sectionwise choice of a set is made.

\paragraph{4. Extend by a dummy phase.}
On $Z$ retain $\psi_y(k)$.  Off $Z$ replace it by the phase $\Psi(k)$ from
\cref{patch:dummy-phase}.  Denote the result by $\psi_y^*(k)$.
Choose its threshold $M$ larger than the frequency bound returned in Step 5, and its exceptional
relative measure smaller than one quarter of the power mass used in Step 6.
These choices are legitimate: Step 5's constants do not depend on phase values.
The modified phase is not asserted to be a cube derivative off $Z$.
For every $k\in K$,
\[
 \int\left|\int G_{y,b,c}(z)\ee{\psi_y^*(k)z}dz\right|dy
                    \geq\delta^{O(1)}N^{D_m}.
\]

\paragraph{5. Apply precisely $\mathsf{MA}(m,l)$.}
Choose a measurable unit phase in $y$ to turn the inner coefficient into its modulus.
Expand $G$, distribute the paired derivative over the factors, and set the new spatial
variable equal to $x+P_m(t)e_m$.  The fixed $\eta$ phase cancels because $r\geq1$.
The resulting integral is
\[
 \int f'_0(x)\mathbb E_t\prod_{i<m}D_{c,b}^{r}f_i(x-P_i(t)e_i)
           \ee{\sum_{i=m}^l\zeta'_i(x)P_i(t)}\,dx,
\]
where
\[
 \zeta'_m(y,z)=-\psi_y^*(k),\qquad
 \zeta'_i(x)=\sum_\omega(-1)^{|\omega|}\xi_i(x+\textstyle\sum_j b_j^\omega e_m)
       \quad(i>m),
\]
and $f'_0$ is the product of the chosen $y$ phase, $D_{c,b}^{r}f_0(x)$, and
$\ee{\psi_y^*(k)x_m}$.  It is $1$-bounded and supported in a budgeted box.
Fix a $z$ section carrying a power-large correlation.  This is exactly
\eqref{patch:ma-correlation}, in $m-1$ variables, with the same $(m,l)$ and zero fixed higher
phase.  The assumed $\mathsf{MA}(m,l)$ yields a power-large set of $y$ on which
$|\psi_y^*(k)|\leq u_{C_6}(\delta)N^{-d_m}$.
Take as this set the full threshold set.  It is jointly measurable in $(y,k)$, whether or
not the intermediate selected $z$ depends measurably on $k$.

\paragraph{6. Recover the original cube phase.}
Choose the dummy threshold in Step 4 greater than $u_{C_6}(\delta)N^{-d_m}$.
Outside its small exceptional set in $k$, every low-frequency point just obtained lies in
$Z$, where $\psi_y^*=\psi_y$.
Removing the exceptional set and applying Fubini/popularity again leaves power-many $y$,
each with power-many $k=(a,b,c)$ for which the original alternating sum is small and all
corners belong to $E_y$.  All losses are bounded by the explicit discard rule in
\cref{patch:budgets}.

\paragraph{7. Quantize and partition the nonzero vertices.}
Partition the interval of possible $\psi_y(k)$ values into cells of width $\tau N^{-d_m}$.
Choose a common cell carrying power mass, with center $\zeta$; the number of cells is
power-controlled.  Choose $\tau>0$ small enough that
$2\pi(\sup_{z\in J-[0,rH]}|z|)\tau N^{-d_m}$ is at most a quarter of the normalized Fourier
coefficient lower bound in Step 1.  This choice is made before the loss from selecting a cell.
For $\Omega_i=\{\omega:\omega_{<i}=0,\omega_i=1\}$ define
\[
 \alpha_1=\zeta-\sum_{\omega\in\Omega_1}(-1)^{|\omega|}\varphi_y(a,b^\omega),
 \qquad
 \alpha_i=-\sum_{\omega\in\Omega_i}(-1)^{|\omega|}\varphi_y(a,b^\omega)\ (i>1).
\]
Each $\alpha_i$ is independent of $b_i$, and on the retained set
\[
 \left|\varphi_y(a,b)-\sum_i\alpha_i(y,a,b,c)\right|
       =|\psi_y(a,b,c)-\zeta|\leq\tau N^{-d_m}.
\]
This is an exact partition of all nonzero Boolean vertices, not an asserted separation of
variables.

\paragraph{8. Remove the phases and integrate.}
For each good $y$, fix one $c$ for which the retained $(a,b)$ section has power-large measure.
On it, replacing the Fourier phase in Step 1 by $\sum_i\alpha_i$ loses at most a quarter of
its lower bound, by the choice of $\tau$ and $|\ee u-\ee v|\leq2\pi|u-v|$.
Enlarge the average to all $(a,b)$ only after taking the squared modulus, which is nonnegative.
Apply \cref{patch:missing-phase}.  Since $r+1=s-1$, it gives a power lower bound for the
order-$(s-1)$ norm of $F_y$ at the same $H$.
There is no need to choose $c$ measurably: the resulting norm inequality is an intrinsic
measurable statement about $F_y$ and holds for every good $y$.
Integrate over the already measurable good set of $y$.  Section normalization ratios,
shift/support ratios, and the fixed roots involved are all handled by \cref{patch:budgets}.

\paragraph{Explicit thresholds for the eight steps.}
Here are numerical choices for every popularity and discretization in this proof. They also
verify that the output budget is uniform. Normalize the transverse envelope $Y$ to a probability
space and the section interval to length $S$, so the full spatial normalizer is $|Y|S$.
Enlarge these two envelopes once at the outset to contain the displayed translates; all
length ratios are bounded above and below by the input budget. Let $0<a\leq1$ be a
fixed lower bound for the resulting order-$s$ power. Put $r=s-2$ and $T=(S+rH)/S$.
Expanding the first $r$ pairs gives an average of $Q_{2,H,S}(D^r_{a,b}F_y)$ at least $a$,
with integrand between zero and one. The set where this integrand is at least $a/2$ has
relative measure at least $a/2$. The set $Y_1$ where its $(a,b)$ section has relative
measure at least $a/4$ has relative measure at least $a/4$.

On those sections rational Fourier selection gives a coefficient threshold
\[
 b=\min\left\{\frac12,\frac{H}{2S}\sqrt{a/2}\right\}>0.
\]
For every $y\in Y_1$ the quantity $\theta$ in the dual--difference lemma is at least
$\theta_0=(a/4)b^2$. Define
\[
 d=\theta_0^{2^r}T^{-2(2^r-1)},\qquad z_0=ad/8.
\]
After dual--difference interchange, let $Z$ be the set on which all corners lie in $E_y$
and the squared normalized coefficient of $G$ is at least $d/2$. That squared coefficient
is at most one: the zero vertex supports $G$ in an interval of length $S$.
Thus every $y\in Y_1$ has a $Z$ section of relative measure at least $d/2$, and the joint
relative measure of $Z$ is at least $z_0$. The set $K$ of shifts for which the $y$ section
has relative measure at least $z_0/2$ has relative measure at least $z_0/2$.
For each $k\in K$, the absolute-coefficient integral in Step 4 is at least
\[
 (z_0/2)\sqrt{d/2}\,|Y|S.
\]
This is the uniform lower threshold fed into Step 5, after conversion to $N^{D_m}$.
To select a $z$ section there, divide half that threshold by the length of the common
$z$ envelope. A section at least this large exists, by the integral triangle inequality.
The assumed major-arc theorem therefore returns a common positive relative lower bound
$\rho$ for the set of $y$ with $|\psi_y^*(k)|\leq B$, where $B=u_A(\delta)N^{-d_m}$.
Decrease $\rho$ to at most one if necessary. Both $\rho$ and $B$ depend only on the budgets
already specified, not on the dummy values.

Choose the dummy threshold greater than $B$ and the exceptional relative measure at most
$z_0/4$. At least $z_0/4$ of the shifts in $K$ remain. For each such shift the returned
low-frequency set has relative $y$ measure at least $\rho$ and lies in $Z$. Consequently
the joint relative measure of original useful small-phase points is at least
$w=z_0\rho/4$.
Let $Z_*\geq N^{d_m}$ also bound $|z|$ on $J-[0,rH]$, and choose frequency-cell width
\[
 \Delta=b/(8\pi Z_*).
\]
Changing a frequency by at most $\Delta$ changes the normalized coefficient in Step 1 by
at most $2\pi Z_*\Delta=b/4$. The interval $[-B,B]$ has a partition into
$q\leq1+\lceil2B/\Delta\rceil$ cells of that width. Some cell carries joint mass at least
$m_0=w/q$. Its set of $y$ sections with relative mass at least $m_0/2$ has relative measure
at least $m_0/2$. For every such $y$, Fubini gives a $c$ section whose relative $(a,b)$ measure
is at least $m_0/4$. There the perturbed coefficient is at least $b/2$. In Step 8 the
quantity $\theta$ of the missing-phase lemma is thus at least $m_0b^2/16$. The final
normalized lower bound is at least
\[
 \frac{m_0}{2}\,
 \frac{(m_0b^2/16)^{2^r}}{T^{2(2^r-1)}(1+H/S)}.
\]
Every factor is positive and is obtained by a fixed number of the budget operations.
The number of degree-lowering iterations is fixed before these choices. This proves the
promised quantitative assertion and specifies the selections without relying on a
uniform finite set of frequencies or a measurable choice of arbitrary sections.
\end{proof}
\paragraph{Dependency warning built into the implementation.}
\pid{conditionalDegreeLowering} is a theorem with an internal premise
\pid{MajorArcProperty m l}; it is proved \emph{before} the major-arc induction.
The next subsection supplies that premise by induction.  The final every-input theorem
will have no such premise.


\subsection{The major-arc induction is noncircular}\label{patch:major-induction}
\begin{theorem}[Real major-arc property]\label{patch:major-arc}
For every increasing degree list contained in $\{1,2,3\}$ and every $1\leq m\leq l\leq k$,
the property $\mathsf{MA}(m,l)$ defined in \eqref{patch:ma-correlation}--\eqref{patch:ma-set}
holds.  All output budgets depend only on the input budget and the degree list.
\end{theorem}
\begin{proof}
Fix $k,l$ and induct on $m=1,\ldots,l$, with the induction hypothesis quantified over
\emph{every} support/coefficient/correlation budget.

\paragraph{Base $m=1$.}
The space $\mathbb R^0$ has one point and measure one.  Since $|g_0(0)|\leq1$, a lower
bound for \eqref{patch:ma-correlation} gives the same lower bound for the scalar exponential
average.  Apply \cref{patch:triangular}.  The canonical set in \eqref{patch:ma-set} is the
whole point space after increasing its budget, and its measure is one.

\paragraph{Induction step: replace the last active input by its exact adjoint.}
Suppose $\mathsf{MA}(m,l)$ is proved, and consider the correlation for $\mathsf{MA}(m+1,l)$:
\[
 \Gamma=\int g_0(x)\mathbb E_t\prod_{i\leq m}g_i(x-P_i(t)e_i)\,
          \ee{p_x(t)}\,dx,\qquad
 p_x(t)=\sum_{i=m+1}^l\zeta_i(x)P_i(t)+\sum_{i>l}\theta_iP_i(t).
\]
Here $x\in\mathbb R^m$.  Define, with all conjugations shown,
\refstepcounter{equation}
\[
 \label{patch:ma-adjoint}
 F(y)=\mathbb E_t\overline{g_0(y+P_m(t)e_m)}
     \prod_{i<m}\overline{g_i(y+P_m(t)e_m-P_i(t)e_i)}\,
               \ee{-p_{y+P_m(t)e_m}(t)}.
 \tag{\theequation}
\]
Translation in $x$ gives $\Gamma=\langle g_m,F\rangle$.
If the $g_i$ have squared $L^2$ norms bounded by $V$, then
\[
 \|F\|_2^2\geq |\Gamma|^2/V.
\]
The function $F$ is $1$-bounded and is supported in a fixed enlarged box, because its first
factor is a translate of $g_0$.  Replacing $g_m$ by $F$ in $\Gamma$ produces exactly
$\langle F,F\rangle=\|F\|_2^2$.  Thus the replaced correlation remains power-large in
units $N^{D_m}$.

\paragraph{Obtain uniformity, using only the smaller major-arc property.}
The phase in the replaced correlation is still $p_x(t)$, a polynomial in $t$ of degree at
most three with measurable coefficients.  Apply \cref{patch:highest-control}, whose proof
is independent of every major-arc assertion.  It gives power-large order-$S_0$ uniformity
of $F$ in direction $e_m$, at a radius $H$ comparable to $N^{d_m}$ up to powers.
The order $S_0$ is bounded independently of all functions and parameters; pad it to a fixed
order if necessary using \eqref{patch:raise-order}.
The function \eqref{patch:ma-adjoint} has exactly the form in
\cref{patch:conditional-degree}, with conjugated $g_i$, mappings $\xi_i=-\zeta_i$, and
constants $\eta_i=-\theta_i$.
Use the induction hypothesis $\mathsf{MA}(m,l)$ in that conditional lemma, and lower the
order $S_0-2$ times, at the same $H$.  If $S_0=2$ there is no lowering step.
This obtains a power-large directional $U^2$ power of $F$.

\paragraph{Fourier-select and expose the smaller pattern.}
By Fubini, popularity, and \cref{patch:u2-fourier-selection}, there is a measurable mapping
$\lambda:\mathbb R^{m-1}\to\mathbb R$ such that
\[
 \int\left|\int F(y,z)\ee{\lambda(y)z}\,dz\right|dy
                      \geq\delta^{O(1)}N^{D_m}.
\]
Take a measurable scalar $c(y)$ of modulus at most one that turns the inner coefficient
into its modulus.  Expand $F$ and put the new variable $x$ equal to the old variable plus
$P_m(t)e_m$.  The resulting power-large scalar is exactly
\refstepcounter{equation}
\[
 \label{patch:ma-smaller-pattern}
 \int \overline{g_0(x)}c(x_{<m})\ee{\lambda(x_{<m})x_m}
 \mathbb E_t\prod_{i<m}\overline{g_i(x-P_i(t)e_i)}
 \ee{-\lambda(x_{<m})P_m(t)-p_x(t)}\,dx.
 \tag{\theequation}
\]
The $x_m$ integration is restricted by $g_0$ to an interval of power-controlled length
in units $N^{d_m}$.  Each section has absolute value at most a power-controlled multiple
of $N^{D_{m-1}}$.  Since the integral of the absolute values of the sections is at least
the modulus of \eqref{patch:ma-smaller-pattern}, popularity gives a measurable set $Z_m$
of $x_m$ values with measure at least $\delta^{O(1)}N^{d_m}$, on every one of which the
section correlation is at least $\delta^{O(1)}N^{D_{m-1}}$.
Use a fixed threshold for all these sections; do not feed their individual sizes into the
induction hypothesis.

For each $z\in Z_m$, the section of \eqref{patch:ma-smaller-pattern} is an
$\mathsf{MA}(m,l)$ correlation, with variable phases
\[
 -\lambda(y),\quad -\zeta_{m+1}(y,z),\ldots,-\zeta_l(y,z)
\]
and fixed higher phases $-\theta_i$.
Apply the induction hypothesis with one common budget.  Drop the nonnegative term
$N^{d_m}|\lambda(y)|$ from its conclusion.  It follows that the section of the canonical set
\[
 Y=\left\{(y,z)\text{ in the enlarged box}:\
       \sum_{i=m+1}^lN^{d_i}|\zeta_i(y,z)|+
       \sum_{i>l}N^{d_i}|\theta_i|\leq u_C(\delta)\right\}
\]
has measure at least $\delta^{O(1)}N^{D_{m-1}}$ for each $z\in Z_m$.
This set is jointly measurable, with no sectionwise set-selection issue.
Fubini yields $|Y|\geq\delta^{O(1)}N^{D_m}$, which is $\mathsf{MA}(m+1,l)$.

All losses in a step are products, fixed powers, or finite-order applications of the budget
rules.  All statements used in this induction hold for every $N>0$ in the real case.
There are at most $k-1\leq2$ induction steps, although each step invokes the
previously bounded number of degree-lowering iterations.  This proves the claimed uniform
structural budgets.
\end{proof}

The dependency at stage $m$ is precisely
\[
 \mathsf{MA}(m,l)
 \ \Longrightarrow\ \mathsf{ConditionalDL}(m,l)
 \ \Longrightarrow\ \mathsf{MA}(m+1,l),
\]
with highest-active uniformity proved before this induction.  In particular, neither the
conclusion at $m+1$ nor the final inverse theorem is used as an induction hypothesis.

\begin{corollary}[Unconditional degree lowering for structured adjoints]
\label{patch:structured-degree}
Let $2\leq l\leq k\leq3$ and fix admissible polynomials and bounded, budget-supported inputs.
Let $F_l$ be the $l$-th adjoint of the modulated average with active inputs $1,\ldots,l$
and a fixed phase $\ee{\sum_{i>l}\zeta_iP_i(t)}$.  If its directional order-$s$ norm,
$s\geq3$, is power-large at $H\asymp\delta^{O(1)}N^{d_l}$, then its order-$(s-1)$ norm
at the same $H$ is power-large, with uniform budgets.
\end{corollary}
\begin{proof}
Keep the passive coordinates $x_{>l}$ fixed.  The active $l$-variable function is exactly
the function in \cref{patch:conditional-degree} with $m=l$ and fixed higher phase,
up to the explicitly harmless conjugations.  Popularity first restricts to passive sections
with power-large order-$s$ power.  Apply that lemma using the now proved
$\mathsf{MA}(l,l)$, with the same budgets on all selected sections, and integrate the
resulting nonnegative order-$(s-1)$ powers.  The passive-section volume has its exact factor
$N^{D-D_l}$; support ratios are handled by \cref{patch:budgets}.
\end{proof}
This corollary is the structured statement corresponding to \cite[Lemma 4.70]{KoszEtAl}.
It makes no assertion that a generic function with large high-order uniformity has large
low-frequency energy or admits degree lowering.
The major-arc induction hypothesis must quantify over budgets; specializing its budget
before the induction would prevent its use on the enlarged supports.


\subsection{Remove the higher inputs, then prove lowest-input energy}
\label{patch:lowest-inverse}
For this subsection write a modulated correlation without inner-product notation as
\[
 \Gamma_l(g_0;f_1,\ldots,f_l;\zeta_{>l})
 =\int g_0(x)\mathbb E_t\prod_{i\leq l}f_i(x-P_i(t)e_i)
                   \ee{\sum_{i>l}\zeta_iP_i(t)}\,dx.
\]
The output factor is $g_0$, not its conjugate.  This convention eliminates implicit
conjugations in the identities below; the original four-linear form corresponds to
$g_0=\overline{f_0}$.

\begin{lemma}[Backward elimination of higher inputs]\label{patch:remove-high-inputs}
For an increasing degree list in $\{1,2,3\}$, a power-large unmodulated correlation
$\Gamma_k(g_0;f_1,\ldots,f_k)$ produces a fixed frequency vector $\zeta_{>1}$ and a
measurable $g'_0$ with $|g'_0|\leq|g_0|$ such that
\[
 |\Gamma_1(g'_0;f_1;\zeta_{>1})|\geq\delta^{O(1)}N^D.
\]
All inputs are $1$-bounded with budgeted support and admissible polynomials.
\end{lemma}
\begin{proof}
Induct backwards on $l=k,k-1,\ldots,2$ while retaining the original $f_1,\ldots,f_l$ and
one fixed vector $\zeta_{>l}$.  The base is the assumed correlation.
Suppose $\Gamma_l$ is power-large.  Define
\[
 F_l(y)=\mathbb E_t\overline{g_{0,l}(y+P_l(t)e_l)}
          \prod_{i<l}\overline{f_i(y+P_l(t)e_l-P_i(t)e_i)}
                           \ee{-\sum_{i>l}\zeta_iP_i(t)}.
\]
Then $\Gamma_l=\langle f_l,F_l\rangle$.  As in \eqref{patch:ma-adjoint}, replacing $f_l$
by $F_l$ gives a power-large correlation equal to $\|F_l\|_2^2$.
Apply \cref{patch:highest-control}, followed by \cref{patch:structured-degree} a bounded
number of times, to obtain power-large directional $U^2$ of $F_l$ at a radius comparable
to $N^{d_l}$ up to powers.

\paragraph{Measurable section frequencies.}
Let $y=x_{<l}$, $z=x_l$, and $w=x_{>l}$.
By \cref{patch:u2-fourier-selection} and popularity there is a measurable set
$Z\subseteq\{(y,w)\}$ of measure at least $\delta^{O(1)}N^{D-d_l}$ and a measurable
frequency $\lambda(y,w)$ for which
\refstepcounter{equation}
\[
 \label{patch:high-section-coeff}
 \left|\int F_l(y,z,w)\ee{\lambda(y,w)z}\,dz\right|
                                  \geq\delta^{O(1)}N^{d_l}\qquad((y,w)\in Z).
 \tag{\theequation}
\]
Popularity in $w$ gives a measurable set $W$ of measure at least
$\delta^{O(1)}N^{D-D_l}$ such that each $Z_w$ has measure at least
$\delta^{O(1)}N^{D_{l-1}}$.
Extend $\lambda$ off $Z$ by a fixed real number $\lambda_{\rm dummy}$ whose modulus is
larger than the major-arc bound used in the next paragraph.  This is permitted because that
bound is independent of the input phase values.  No Fourier-coefficient assertion is made
for the extended values off $Z$.

\paragraph{Use $\mathsf{MA}(l,l)$ to keep the selected frequencies small.}
For each $w\in W$, the integral over $y$ of the absolute coefficient with this extended
frequency is power-large.  Dualize by a measurable phase in $y$, expand $F_l$, and translate
by $P_l(t)e_l$.  This gives a power-large integral in $(y,z)$ with the $l-1$ active inputs
$\overline{f_i}$, the variable coefficient $-\lambda(y,w)$ of $P_l(t)$, and fixed coefficients
$-\zeta_i$ for $i>l$.
Select a $z$ section carrying a power-large correlation and apply the already proved
$\mathsf{MA}(l,l)$.  Its output budget is the same for every $w$.
In particular, the canonical set of $y$ with
\[
 N^{d_l}|\lambda(y,w)|+\sum_{i>l}N^{d_i}|\zeta_i|\leq u_C(\delta)
\]
is power-large.  Choose $|\lambda_{\rm dummy}|>u_C(\delta)N^{-d_l}$.
Every point of this canonical set must lie in $Z_w$, so the original estimate
\eqref{patch:high-section-coeff} holds there.  The full set of these $(y,w)$ is jointly
measurable, is power-large, and has $|\lambda(y,w)|\leq u_C(\delta)N^{-d_l}$.
This dummy-value step prevents the major-arc theorem from returning unrelated points on
which no useful Fourier coefficient was obtained.

\paragraph{Make the frequency constant.}
Partition that bounded frequency interval into cells of width $\tau N^{-d_l}$.
Choose $\tau>0$ so that changing $\lambda$ by at most $\tau N^{-d_l}$ changes the coefficient
in \eqref{patch:high-section-coeff} by at most one quarter of its lower bound.
Indeed $F_l$ is $1$-bounded and its $z$ support has length and absolute location bounded
by $u_c(\delta)N^{d_l}$; integrate
$|\ee{\lambda z}-\ee{\lambda_0 z}|\leq2\pi|\lambda-\lambda_0||z|$.
The chosen $\tau$ is a power of $\delta$ before the cell selection.
There are power-many cells, so a common center $\lambda_0$ works on a measurable set of
$(y,w)$ of power-large measure.  The coefficient at $\lambda_0$ is still pointwise power-large
on that set.  Choose a measurable $c(y,w)$, zero off that set and of modulus one on it,
such that
\[
 \int c(y,w)\int F_l(y,z,w)\ee{\lambda_0z}\,dz\,dy\,dw
\]
is nonnegative and power-large.
Expand $F_l$, translate by $P_l(t)e_l$, and conjugate the entire scalar integral.
The result is exactly $\Gamma_{l-1}$, with
\[
 g_{0,l-1}(x)=g_{0,l}(x)\overline{c(x_{\{l\}^c})}\ee{-\lambda_0x_l},
 \qquad \zeta_l=\lambda_0,
\]
and with the earlier fixed $\zeta_{>l}$ unchanged.  Thus
$|g_{0,l-1}|\leq|g_{0,l}|$, the support does not enlarge, and the remaining lower inputs are
still the original $f_i$, not modified products.
There are at most two backward steps.  All budget choices are explicit applications of
\cref{patch:budgets}, so the iteration yields the stated conclusion.
\end{proof}
The inequality $|g'_0|\leq|g_0|$ is the form needed here.  The paper states a modulus-equality
version in Lemma 4.71; no such stronger equality is needed or asserted in this patch.

\begin{theorem}[Lowest-input positive Fourier energy]\label{patch:lowest-energy}
Fix the increasing degree list and an input budget.  There is a structural integer $C$ such
that a power-large correlation $\Gamma_k$ as above, for $N\geq u_C(\delta)$, implies
\[
 \left\langle f_1,P_{L}^{(1)}f_1\right\rangle
       =\int\eta(\xi_1/L)|\widehat f_1(\xi)|^2\,d\xi
       \geq\ell_C(\delta)N^D,
       \qquad L=u_C(\delta)N^{-d_1}.
\]
The energy is real and nonnegative.  The theorem holds for every increasing sublist of
$\{1,2,3\}$, including lists of length one whose sole degree is two or three.
\end{theorem}
\begin{proof}
For $k\geq2$, apply \cref{patch:remove-high-inputs}.  For $k=1$ begin with the original
correlation.  Set
\[
 T f_1(x)=\mathbb E_t f_1(x-P_1(t)e_1)\ee{\sum_{i>1}\zeta_iP_i(t)}.
\]
Cauchy--Schwarz with the supported output factor gives
$\|Tf_1\|_2^2\geq\delta^{O(1)}N^D$.
\Cref{new:average-multiplier,new:integral-plancherel} give
\[
 \|Tf_1\|_2^2=\int |m(\xi_1)|^2|\widehat f_1(\xi)|^2\,d\xi,
 \qquad m(v)=\mathbb E_t\ee{-vP_1(t)+\sum_{i>1}\zeta_iP_i(t)}.
\]
We have $|m|\leq1$ and $\|f_1\|_2^2\leq u_c(\delta)N^D$ for a known budget.
Write the already obtained bounds as $\|Tf_1\|_2^2\geq\kappa N^D$
and $\|f_1\|_2^2\leq V_0N^D$, where $\kappa>0$ is power-large and $V_0>0$ is
power-bounded. Choose
\[
 \varepsilon=\min\{1/2,\sqrt{\kappa/(4V_0)}\}>0.
\]
Then $\varepsilon^2\|f_1\|_2^2\leq\kappa N^D/4$.
By \cref{patch:triangular}, $|m(v)|\geq\varepsilon$ implies
$|v|\leq R=u_{C'}(\delta)N^{-d_1}$.
Thus the Fourier energy in $|\xi_1|\leq R$ is power-large.
Choose $L\geq4R$.  The cutoff equals one on this band and is nonnegative elsewhere, so
\[
 \langle f_1,P_L^{(1)}f_1\rangle\geq
          \int_{|\xi_1|\leq R}|\widehat f_1(\xi)|^2d\xi
                              \geq\delta^{O(1)}N^D.
\]
Choose a single integer $C$ dominating the threshold, radius coefficient and exponent,
and lower-bound coefficient and exponent, so that $L=u_C(\delta)N^{-d_1}$ still has its
plateau on the same band.  This proves the exact statement.
\end{proof}
This last step uses the source's modulated-linear-average and scalar-oscillation argument,
but full Plancherel avoids an unnecessary final sectionwise choice of Fourier data.

\subsection{General selected indices: quantitative freezing and slicing}
\label{patch:all-indices}
The next theorem is the independent substitute for the source's Theorem 4.30 needed by the
public monomial wrappers.  Its proof is given here rather than taken as an axiom.

\begin{theorem}[Every-input positive energy, independently proved]
\label{patch:energy-core}
Fix an increasing list $1\leq d_1<\cdots<d_k\leq3$ and an integer input budget $c\geq1$.
There exists an integer $C\geq c$ such that the following holds for $0<\delta\leq1$ and
$N\geq u_C(\delta)$.  Suppose that all $P_i$ are admissible with budget $c$, the functions
$g_0,f_1,\ldots,f_k$ are $1$-bounded and supported in the budget-$c$ box, and
\[
 |\Gamma_k(g_0;f_1,\ldots,f_k)|\geq\ell_c(\delta)N^D.
\]
Then, for every $j\in\{1,\ldots,k\}$,
\refstepcounter{equation}
\[
 \label{patch:all-energy}
 \left\langle f_j,P_{L_j}^{(j)}f_j\right\rangle
      \geq\ell_C(\delta)N^D,
       \qquad L_j=u_C(\delta)N^{-d_j}.
 \tag{\theequation}
\]
\end{theorem}
\begin{proof}
Induct on $k$, with the statement quantified over all increasing degree lists and all input
budgets.  For $k=1$, or for $j=1$ at any $k$, use \cref{patch:lowest-energy}.
Suppose $k\geq2$ and $j>1$.

\paragraph{1. Replace the first input, and smooth that adjoint.}
Define its exact adjoint
\[
 F_1(y)=\mathbb E_t\overline{g_0(y+P_1(t)e_1)}
                  \prod_{i>1}\overline{f_i(y+P_1(t)e_1-P_i(t)e_i)}.
\]
The original scalar is $\langle f_1,F_1\rangle$.  Its size gives
$\|F_1\|_2^2\geq\delta^{O(1)}N^D$, and replacing $f_1$ by $F_1$ produces this positive
correlation.  All its inputs, including $F_1$, are $1$-bounded and supported in a budgeted
enlargement.  Apply \cref{patch:lowest-energy} to this new correlation to obtain
\[
 \langle F_1,P_{L_1}^{(1)}F_1\rangle\geq\kappa N^D,
 \qquad L_1=\delta^{-O(1)}N^{-d_1},\quad\kappa=\delta^{O(1)}.
\]
Set $G=P_{L_1}^{(1)}F_1$.  The adjoint identity gives the same lower bound for the
correlation with first input $G$ and other inputs the original $f_2,\ldots,f_k$.
The scalar is real by self-adjointness, but only its modulus will be used.

Let $M=\max(1,\|\check\eta\|_1)$ and $M_1=\| (\check\eta)'\|_1$.
The kernel estimate in \cref{new:kernel-facts} gives
\[
 \|G\|_\infty\leq M,\qquad
 |G(x+ve_1)-G(x)|\leq M_1L_1|v|.
\]
This holds on every bounded Borel section directly from the $L^1$ translation bound for the
smooth kernel. No differentiability of $F_1$ or of a merely measurable input is assumed.
The output factor $g_0$ retains finite support even though $G$ itself need not have compact
support in coordinate one.

\paragraph{2. Freeze on one of finitely many equal intervals.}
Admissibility bounds $|P_1'(t)|\leq u_{c'}(\delta)N^{d_1-1}$ on $[0,N]$.
Partition $[0,N]$ into $K$ equal intervals $[a,a+n]$ with $n=N/K$.
On such an interval,
\[
 |G(x-P_1(t)e_1)-G(x-P_1(a)e_1)|
 \leq M_1L_1u_{c'}(\delta)N^{d_1-1}n
 \leq u_{c''}(\delta)/K.
\]
After integrating against $g_0$ and the other bounded factors, the total normalized
correlation error is at most $u_{c'''}(\delta)N^D/K$.
Choose
\[
 K=\max\{1,\lceil 4u_{c'''}(\delta)/\kappa\rceil\}.
\]
This is power-bounded above, and ensures error at most $\kappa N^D/4$.
The full average is the average of these $K$ normalized local averages.  Hence at least one
frozen interval has scalar modulus at least $\kappa N^D/2$.
Choose the first such interval in this finite family.  No discarded remainder interval is needed.

\paragraph{3. Expose the lower-dimensional polynomial family.}
For that $a$, put
\[
 Q_i(u)=P_i(a+u)-P_i(a),\qquad
 f_i^a(x)=f_i(x-P_i(a)e_i)\quad(i\geq2),
\]
\[
 g_0^a(x)=M^{-1}g_0(x)G(x-P_1(a)e_1).
\]
The frozen normalized correlation, divided by $M$, is exactly
\refstepcounter{equation}
\[
 \label{patch:frozen-correlation}
 \int g_0^a(x)\mathbb E_{u\in[0,n]}
               \prod_{i=2}^k f_i^a(x-Q_i(u)e_i)\,dx.
 \tag{\theequation}
\]
It is power-large in units $N^D$, and all its displayed functions are $1$-bounded.
The $Q_i$ have degrees $d_i$ and the same leading coefficients as $P_i$.
By the binomial formula and $0\leq a\leq N$, their degree-$r$ coefficients are bounded by
$u_{c'}(\delta)N^{d_i-r}$, hence by
$u_{c'}(\delta)K^{d_i-r}n^{d_i-r}$.
Since $K$ is power-bounded, these are admissible coefficients at scale $n$ with one new
structural budget.  The translated supports are contained in boxes with sides bounded by
$u_{c'}(\delta)N^{d_i}=u_{c'}(\delta)K^{d_i}n^{d_i}$ and therefore satisfy the same kind
of enlarged-support hypothesis at scale $n$.

\paragraph{4. Obtain a measurable set of large sections.}
For $v=x_1$, let $B(v)$ be the integral in the remaining coordinates in
\eqref{patch:frozen-correlation}.  It is measurable by Fubini, vanishes unless $v$ lies in
the original coordinate-one support interval, and
$|B(v)|\leq u_{c'}(\delta)N^{D-d_1}$.
Since $\int|B(v)|dv\geq|\int B(v)dv|\geq\delta^{O(1)}N^D$, popularity produces a
canonical measurable set $X_1$ of measure at least $\delta^{O(1)}N^{d_1}$ on which
\[
 |B(v)|\geq\delta^{O(1)}N^{D-d_1}\geq\delta^{O(1)}n^{D-d_1}.
\]
All hypotheses of the $(k-1)$-dimensional induction theorem now hold with the increasing
list $(d_2,\ldots,d_k)$ and one common input budget, independent of $v\in X_1$.
This includes lists starting with degree two or three, not only lists starting with one.
Because $n=N/K$ and $K\leq u_{c'}(\delta)$, choosing the original threshold $N$ larger
by the product of the two known budget bounds ensures the required lower threshold for $n$.

\paragraph{5. Apply the induction hypothesis and integrate positive energies.}
The induction hypothesis gives, for every $v\in X_1$, a lower bound
$\delta^{O(1)}n^{D-d_1}$ for the positive energy of the section of $f_j^a$ at the same
radius $L'=u_{C'}(\delta)n^{-d_j}$.  Translation by $P_j(a)e_j$ does not change its
coordinate-one section and preserves the directional Fourier energy in the remaining
coordinates.  Thus this is also the section energy of the original $f_j$.
Since $n=N/K$ and $K$ is power-bounded, the lower bound is
$\delta^{O(1)}N^{D-d_1}$.
Integrate over $X_1$ and use nonnegativity on its complement.  This gives
\[
 \langle f_j,P_{L'}^{(j)}f_j\rangle\geq\delta^{O(1)}N^D,
 \qquad L'\leq u_{C''}(\delta)N^{-d_j}.
\]
The radius is common to all sections; no varying-radius integral is being identified with
one projection.
Finally choose $L_j=u_C(\delta)N^{-d_j}\geq2L'$.  On the support of
$\eta(\xi_j/L')$, the larger cutoff equals one, so
$\eta(\xi_j/L_j)\geq\eta(\xi_j/L')$ pointwise.
The positive energy therefore only increases.  A common maximum of the finitely many
budgets handles all $j$, thresholds, and the degree-list induction.
This proves \eqref{patch:all-energy}.
\end{proof}

\paragraph{Why this is not the old circular proof.}
\Cref{patch:energy-core} uses the independently proved \cref{patch:lowest-energy}, the
adjoint identities, freezing, and lower-dimensional instances of itself.
It does not use the public \pid{lem:pet-reduction}, the old \pid{lem:degree-lowering-zero},
or the public \pid{thm:real-inverse}.  The only PET premise used in proving
\cref{patch:lowest-energy} was proved explicitly in \cref{patch:pet-producer}.


\section{The monomial inverse theorem and the compact adjoint estimate}
\label{new:compact-section}
Return to $\R^3$ and to the positive translations $t,t^2,t^3$. Write
\[
 B_N(C_0)=[-C_0N,C_0N]\times[-C_0N^2,C_0N^2]\times[-C_0N^3,C_0N^3].
\]
The pairing is linear in its first argument, as in the conventions.

\begin{theorem}[The precise inverse statement used in the adjoint argument]\label{new:monomial-inverse}
For every $C_0\geq1$ there is an integer $A\geq1$ such that, for $0<\delta\leq1$,
$N\geq A\delta^{-A}$, and Borel functions $f_0,f_1,f_2,f_3$ bounded by one and supported
in $B_N(C_0)$, the implication
\[
 |\langle A_N(f_1,f_2,f_3),f_0\rangle|\geq\delta N^6
\]
has the following conclusion, for every $j\in\{1,2,3\}$:
\[
 \left\langle f_j,P_{A\delta^{-A}N^{-j}}^{(j)}f_j\right\rangle
 \geq A^{-1}\delta^A N^6.
\]
This scalar is real and nonnegative. In particular the witness in a formulation asking for
$\langle f_j,P_L^{(j)}h_j\rangle$ can be chosen to be $h_j=f_j$.
\end{theorem}
\begin{proof}
Apply \cref{patch:energy-core} with degrees $(1,2,3)$, polynomials $P_i(t)=-t^i$, and
$g_0=\overline{f_0}$. These signs give $x-P_i(t)e_i=x+t^ie_i$, so its correlation is the
one displayed here, not its adjoint or a different translation convention. Choose an integer
input budget $c\geq\max\{1,\lceil C_0\rceil\}$. Then $|\operatorname{lc}P_i|=1$, the lower
coefficients vanish, and the support conditions hold with that budget. Also
$\delta\geq\ell_c(\delta)$. The output budget $C$ in \cref{patch:energy-core} supplies the
claimed conclusion with $A=C$. Any further increase of $A$ may be made with $A\geq2C$;
then the larger cutoff dominates the old one on its support by
\cref{new:multiplier-facts}, while its lower bound weakens and its scale threshold strengthens.
This permits one common integer for all three indices. Nonnegativity follows from the same
multiplier lemma. No instance of the source's inverse theorem is being assumed in this step.
\end{proof}

\begin{lemma}[Adjoint support and the elementary global bound]\label{new:adjoint-basic}
Let $\mathcal G=A_N^{*j}(g_0,g_b,g_c)$, with $j,b,c$ distinct. If the inputs are bounded
by one and supported in $B_N(C_0)$, then $|\mathcal G|\leq1$ and
$\operatorname{supp}\mathcal G\subseteq B_N(C_0+1)$.
For exponents $r_i\in[1,\infty]$, $i\in\{0,b,c\}$, with
$1/r=1/r_0+1/r_b+1/r_c$ and $1\leq r<\infty$,
\[
 \|\mathcal G\|_r\leq\|g_0\|_{r_0}\|g_b\|_{r_b}\|g_c\|_{r_c},
\]
\[
 \|(I-P_R^{(j)})\mathcal G\|_r
 \leq(1+B_0)\|g_0\|_{r_0}\|g_b\|_{r_b}\|g_c\|_{r_c}.
\]
These estimates hold for all inputs with finite displayed norms.
\end{lemma}
\begin{proof}
The pointwise integrand has modulus at most one. If it is nonzero, its $g_0$ argument
$y-t^je_j$ lies in $B_N(C_0)$; since $0\leq t^j\leq N^j$, the output lies in the asserted
enlargement. The norm bound is H\"older for the three translated inputs at each fixed $t$,
followed by the integral inequality and translation invariance from
\cref{new:integral-inequalities}. The second estimate follows from the convolution bound
$\|P_R^{(j)}\|_{L^r\to L^r}\leq B_0$. Approximation by bounded supported functions, or
Tonelli applied to their absolute values followed by this norm bound, gives the stated
extension and its almost-everywhere interpretation as the concrete parameter integral.
\end{proof}

\begin{theorem}[Compact bounded-input high-frequency estimate]\label{new:compact-highpass}
Fix $C_0\geq1$ and $j$. There are $a\in(0,1)$, $K_0\geq1$, and $C_1\geq1$ such that
for $0<\mu\leq1$, $N\geq K_0\mu^{-K_0}$, and bounded Borel inputs supported in $B_N(C_0)$,
\[
 \left\|(I-P_R^{(j)})A_N^{*j}(g_0,g_b,g_c)\right\|_2
 \leq C_1\mu^a N^3\|g_0\|_\infty\|g_b\|_\infty\|g_c\|_\infty,
 \qquad R=N^{-j}\mu^{-2}.
\]
\end{theorem}
\begin{proof}
Zero input norms give zero output. By homogeneity reduce to inputs bounded by one, using
Borel representatives with pointwise bounds. Put $\mathcal G=A_N^{*j}(g)$ and
$H=(I-P_R^{(j)})\mathcal G$. By \cref{new:adjoint-basic}, $\mathcal G$ is supported in
$B_N(C_0+1)$, is bounded by one, and belongs to $L^1\cap L^2$. The same membership for
$H$ follows from the convolution estimates. Let
\[
 M=1+B_0,\qquad B'=B_N(C_0+2),\qquad h=M^{-1}\one_{B'}H.
\]
Then $|h|\leq1$, and $h$ has the same fixed support envelope as the three other inputs.
Since $\mathcal G$ is zero off $B'$, the multiplier identities give
\[
 \langle h,\mathcal G\rangle
 =M^{-1}\int(1-\eta(\xi_j/R))|\widehat{\mathcal G}(\xi)|^2\,d\xi
 \geq M^{-1}\|H\|_2^2.
\]
The inequality uses $0\leq1-\eta\leq1$, so $(1-\eta)^2\leq1-\eta$. In particular this
correlation is real and nonnegative.

Take the integer $A$ from \cref{new:monomial-inverse} with support parameter $C_0+2$ and set
\[
 \delta=\mu^{1/(2A)},\qquad L=A\delta^{-A}N^{-j}=A\mu^{-1/2}N^{-j}.
\]
Choose $K_0\geq A$ and $K_0\geq1$. Then $N\geq K_0\mu^{-K_0}$ implies
$N\geq A\mu^{-1/2}=A\delta^{-A}$. Suppose, towards a contradiction, that
\[
 \|H\|_2^2>M\delta N^6.
\]
The adjoint identity in \cref{new:improving}, applied with the $j$-th input equal to $h$,
turns $\langle h,\mathcal G\rangle$ into the original four-function correlation.
The inverse theorem therefore implies
\[
 \langle h,P_L^{(j)}h\rangle\geq A^{-1}\mu^{1/2}N^6.
\]

For all sufficiently small $\mu$, $2A\mu^{3/2}\leq1$, so $2L\leq R$.
Consequently $P_L^{(j)}H=0$ in $L^2$, by \cref{new:multiplier-facts}. Hence
\[
 P_L^{(j)}h=-M^{-1}P_L^{(j)}(\one_{(B')^c}H).
\]
Outside $B'$, $H=-P_R^{(j)}\mathcal G$. The convolution vanishes outside the passive-coordinate
support of $\mathcal G$. At any remaining point outside $B'$, a nonzero convolution
integrand must have $|u|\geq N^j$: the $j$ coordinate of the point has distance at least
$N^j$ from the support of $\mathcal G$. The pointwise integral triangle inequality and Young's
inequality therefore give
\[
 \|\one_{(B')^c}H\|_2
 \leq\left(\int_{|u|\geq N^j}|K_R(u)|\,du\right)\|\mathcal G\|_2
 \leq B_2^{\rm tail}\mu^4[2(C_0+1)]^{3/2}N^3.
\]
Use the $L^2$ contraction of $P_L^{(j)}$, $\|h\|_2\leq[2(C_0+2)]^{3/2}N^3$, and
Cauchy--Schwarz. There is a fixed finite $C_2$ such that
\[
 0\leq\langle h,P_L^{(j)}h\rangle\leq C_2\mu^4N^6.
\]
Choose $0<\mu_0\leq1$ so that both $2A\mu_0^{3/2}\leq1$ and
$C_2\mu_0^{7/2}\leq(2A)^{-1}$. If $C_2=0$ the latter restriction is unnecessary.
For $\mu\leq\mu_0$ the last two energy bounds contradict each other. Thus
\[
 \|H\|_2\leq\sqrt M\,\mu^{1/(4A)}N^3\qquad(0<\mu\leq\mu_0).
\]
For $\mu_0\leq\mu\leq1$, integral Plancherel gives the crude bound
$\|H\|_2\leq\|\mathcal G\|_2\leq[2(C_0+1)]^{3/2}N^3$; enlarge the constant by
$\mu_0^{-1/(4A)}$. Set $a=1/(4A)$ and take a constant at least one covering both cases.
This proves the theorem. The proof does not assume compact support of $P_R\mathcal G$;
its noncompact part is exactly the tail estimated above.
\end{proof}

\begin{theorem}[A compact decaying point on the $L^2$ exponent hyperplane]
\label{new:compact-decaying-point}
Fix $j$ and use the choice of $b,c$ made in \cref{new:improving-section}. Put
\[
 \rho_0=\frac{10}{61},\qquad\rho_b=\frac{21}{122},\qquad
 \rho_c=\frac{10}{61},\qquad u_i=1/\rho_i.
\]
Thus $\sum_{i\in\{0,b,c\}}\rho_i=1/2$ and all $u_i>2$. There are constants
$a_1\in(0,1)$, $K_0,C_3\geq1$ such that, for inputs supported in $B_N(C_0)$,
\[
 \|(I-P_{N^{-j}\mu^{-2}}^{(j)})A_N^{*j}(g)\|_2
 \leq C_3\mu^{a_1}\prod_{i\in\{0,b,c\}}\|g_i\|_{u_i}
\]
whenever $0<\mu\leq1$ and $N\geq K_0\mu^{-K_0}$.
\end{theorem}
\begin{proof}
For fixed $N,\mu$, the compact bounded-input endpoint has constant $C_1\mu^a N^3$.
The improving endpoint \cref{new:improving} has constant $CN^{-1/20}$ and input exponents
$(6,40/7,6)$. Composing with $I-P_R$ does not enlarge its $L^2$ constant: its multiplier
lies in $[0,1]$, so it is a contraction. Apply \cref{new:interpolation} on the fixed support
box with weight $\theta=60/61$ on the improving endpoint and $1/61$ on the bounded endpoint.
Conjugate the $b,c$ inputs first to make the operator complex multilinear, as allowed there.
The three interpolated reciprocals are exactly the displayed $\rho_i$, and
\[
 3(1-\theta)-\theta/20=3/61-3/61=0.
\]
The frequency gain is $\mu^{a/61}$, so take $a_1=a/61$. Initially the proof applies to
bounded finite-valued inputs with the specified supports. They are dense in each $L^{u_i}$
on that box. The crude global estimate in \cref{new:adjoint-basic}, with output two and
$\sum1/u_i=1/2$, shows that the outputs converge in $L^2$ under these approximations:
expand a difference of two triples as three multilinear differences and use boundedness of
the other factors' norms. The same estimate for absolute inputs identifies the limit with
the concrete parameter integral and its convolution almost everywhere. Passing to the limit
proves the stated compact estimate for all finite-norm inputs.
\end{proof}

\section{Support removal, all exponents, and the exact Lean proposition}
\label{new:global-section}
\subsection{Support removal with an explicit summable off-diagonal estimate}
\begin{theorem}[Global decaying $L^2$ estimate]\label{new:global-decaying-point}
With the exponents and constants of \cref{new:compact-decaying-point}, there are
$C_4,K_0\geq1$ such that
\[
 \|(I-P_{N^{-j}\mu^{-2}}^{(j)})A_N^{*j}(g)\|_2
 \leq C_4\mu^{a_1}\prod_{i\in\{0,b,c\}}\|g_i\|_{u_i}
\]
for all $g_i\in L^{u_i}(\R^3)$, with no support restriction, provided
$0<\mu\leq1$ and $N\geq K_0\mu^{-K_0}$.
\end{theorem}
\begin{proof}
First take bounded inputs of bounded support, so that all sums introduced here are finite.
For $\nu\in\mathbb Z^3$ let
\[
 Q_\nu=D_N(\nu+[0,1)^3),\qquad g_{i,\nu}=\one_{Q_\nu}g_i.
\]
These boxes partition $\R^3$. In the adjoint integral put $z=y-t^je_j$. If $z\in Q_\nu$,
then $z+t^be_b$ lies in $Q_{\nu+\epsilon e_b}$ for exactly one $\epsilon\in\{0,1\}$,
and $z+t^ce_c$ lies in $Q_{\nu+\epsilon'e_c}$ for one $\epsilon'\in\{0,1\}$.
This follows coordinate by coordinate from $0\leq t^i\leq N^i$ and the half-open convention.
Consequently the exact decomposition is
\[
 A_N^{*j}(g)=\sum_{\epsilon,\epsilon'\in\{0,1\}}\sum_{\nu\in\mathbb Z^3}
 A_N^{*j}(g_{0,\nu},g_{b,\nu+\epsilon e_b},g_{c,\nu+\epsilon'e_c}).
\]
Fix one of the four offset pairs, and denote its individual adjoints by $G_\nu$ and its
high-frequency outputs by $H_\nu=(I-P_R^{(j)})G_\nu$.
After the common translation $D_N\nu$, all their input supports are in $B_N(2)$.
Translation covariance of both the average and the convolution, followed by
\cref{new:compact-decaying-point} with $C_0=2$, gives
\[
 \|H_\nu\|_2\leq C\mu^{a_1}a_\nu,
 \qquad
 a_\nu=\|g_{0,\nu}\|_{u_0}\|g_{b,\nu+\epsilon e_b}\|_{u_b}
                       \|g_{c,\nu+\epsilon'e_c}\|_{u_c}.
\]
The crude bound \cref{new:adjoint-basic} gives $\|G_\nu\|_2\leq a_\nu$.
Moreover $G_\nu$ is supported in the same two passive-coordinate intervals as $Q_\nu$, and
in $[\nu_jN^j,(\nu_j+2)N^j)$ in coordinate $j$. The one-coordinate convolution preserves
those passive-coordinate supports exactly.

Let $Q_\omega$ be an output box. Its intersection with $H_\nu$ is empty unless
$\omega_i=\nu_i$ for $i\ne j$. Put $\ell=\omega_j-\nu_j$ in the remaining case.
For $|\ell|\leq3$ use the compact estimate. For $|\ell|>3$, the unprojected $G_\nu$
vanishes on $Q_\omega$ and a nonzero convolution term must have
$|u|\geq(|\ell|-2)N^j$. The same tail-Young calculation as in
\cref{new:compact-highpass}, now with $\|G_\nu\|_2\leq a_\nu$, gives
\[
 \|\one_{Q_\omega}H_\nu\|_2\leq B_2^{\rm tail}\mu^4(|\ell|-2)^{-2}a_\nu.
\]
Define the nonnegative sequence on $\mathbb Z$
\[
 w_\ell=C\mu^{a_1}\quad (|\ell|\leq3).
\]
\[
 w_\ell=B_2^{\rm tail}\mu^4(|\ell|-2)^{-2}\quad (|\ell|>3).
\]
The series $\sum_{n\geq2}n^{-2}$ is finite (compare it to the integral of $x^{-2}$), and
$\mu^4\leq\mu^{a_1}$. Thus $\|w\|_{\ell^1}\leq C'\mu^{a_1}$.
For $b_\omega=\|\one_{Q_\omega}\sum_\nu H_\nu\|_2$, the triangle inequality yields
\[
 b_\omega\leq\sum_{\ell\in\mathbb Z}w_\ell a_{\omega-\ell e_j}.
\]
Take the $\ell^2(\mathbb Z^3)$ norm. The triangle inequality for finite sums, followed by
monotone convergence, proves the discrete Young bound
\[
 \|b\|_{\ell^2}\leq\sum_\ell w_\ell
           \left(\sum_\omega a_{\omega-\ell e_j}^2\right)^{1/2}
 =\|w\|_{\ell^1}\|a\|_{\ell^2}.
\]
There is no overlap loss hidden here: disjointness of the output boxes gives
$\|\sum_\nu H_\nu\|_2=\|b\|_{\ell^2}$. Discrete H\"older, with exponents $u_i/2>1$
and $\sum_i2/u_i=1$, gives
\[
 \|a\|_{\ell^2}\leq
 \prod_{i\in\{0,b,c\}}\left(\sum_\nu\|g_{i,\nu}\|_{u_i}^{u_i}\right)^{1/u_i}
 =\prod_i\|g_i\|_{u_i}.
\]
The fixed offsets merely permute the lattice in their corresponding sums. Sum the four
offset estimates by the $L^2$ triangle inequality. This proves the global bound for the
initial bounded inputs. Approximate arbitrary $L^{u_i}$ inputs by bounded finite-support
ones. As in \cref{new:compact-decaying-point}, the crude global $L^2$ estimate proves
convergence and identifies the limiting operator; the improved bound passes to the limit.

This decomposition is performed on the \emph{inputs}. Cutting the unprojected output into
boxes before applying $I-P_R$ would introduce unestimated boundary terms and is not used.
\end{proof}

\subsection{Explicit convex completion of the exponent range}
\begin{theorem}[The real monomial adjoint estimate at every required exponent]
\label{new:adjoint-all-exponents}
Fix $j\in\{1,2,3\}$ and numbers $\alpha_0,\alpha_1,\alpha_2,\alpha_3\in(0,1)$ with
$\sum_{i=0}^3\alpha_i=1$. There are constants $c\in(0,1)$ and $K,C\geq1$ such that, for
$0<\mu\leq1$, $N\geq K\mu^{-K}$, and $g_i\in L^{1/\alpha_i}(\R^3)$ for $i\ne j$,
\[
 \left\|(I-P_{N^{-j}\mu^{-2}}^{(j)})A_N^{*j}((g_i)_{i\ne j})\right\|_{1/(1-\alpha_j)}
 \leq C\mu^c\prod_{i\ne j}\|g_i\|_{1/\alpha_i}.
\]
There is no compact-support assumption.
\end{theorem}
\begin{proof}
Let $b,c_*$ denote the other two positive indices, to distinguish the index from the gain
constant. Set $\sigma=\alpha_0+\alpha_b+\alpha_{c_*}=1-\alpha_j\in(0,1)$.
Use the fixed reciprocals $\rho_i$ of \cref{new:global-decaying-point}. Choose
\[
 \tau=\frac14\min\left\{1,\ 2(1-\sigma),\
       \min_{i\ne j}\frac{\alpha_i}{\rho_i},\
       \min_{i\ne j}\frac{1-\alpha_i}{1-\rho_i}\right\}>0.
\]
All entries are strictly positive and finite, so $0<\tau\leq1/4<1$.
Put $v_i=(\alpha_i-\tau\rho_i)/(1-\tau)$ for $i\ne j$.
The choice of $\tau$ gives $0<v_i<1$. Indeed,
$\tau\rho_i\leq\alpha_i/4$ makes the numerator positive, and
$\tau(1-\rho_i)\leq(1-\alpha_i)/4<1-\alpha_i$ implies $v_i<1$.
Also
\[
 0<\sum_{i\ne j}v_i=\frac{\sigma-\tau/2}{1-\tau}<1.
\]
Positivity follows from positivity of every $v_i$. The upper inequality is equivalent to
$\tau<2(1-\sigma)$, which is built into the choice. Put $r=1/\sum v_i>1$.
The crude global estimate \cref{new:adjoint-basic} gives a uniform bound from
$\prod_i L^{1/v_i}$ to $L^r$ with constant $1+B_0$.
Interpolate this with \cref{new:global-decaying-point}, assigning weight $\tau$ to the
latter decaying estimate. The exact reciprocal identities are
\[
 \alpha_i=(1-\tau)v_i+\tau\rho_i,
 \qquad \sigma=(1-\tau)/r+\tau/2.
\]
\Cref{new:interpolation} therefore gives the desired input and output exponents and the
gain $\mu^{a_1\tau}$. Set $c=a_1\tau\in(0,1)$, keep the threshold constant $K$ from the
global decaying point (enlarging it to at least one), and enlarge the interpolation constant
to $C\geq1$.

Initially this interpolation is on bounded finite-valued supported functions. Approximate
all three finite-norm inputs by such functions. Since their reciprocals sum to $\sigma<1$,
\cref{new:adjoint-basic} with output $1/\sigma$ bounds each term in the telescoping
multilinear difference and proves norm convergence. Its absolute-value version supplies
an almost-everywhere defined concrete average, and Young's inequality supplies its
convolution. The resulting output is the same limit. Thus the estimate holds for the
stated full $L^p$ domains.
\end{proof}

\begin{corollary}[The exact proposition \texttt{Auto.KoszAdjoint}]\label{new:exact-lean-target}
For every $j$ of type \pid{Fin 3}, the proposition \pid{Auto.KoszAdjoint j} defined in the supplied
source is true. In mathematical indexing, its complete statement is as follows.
For $\alpha_0\in(0,1)$, $\alpha_i\in(0,1)$ for $1\leq i\leq3$, and
$\alpha_0+\sum_{i=1}^3\alpha_i=1$, there are $c\in(0,1)$, $K,C\geq1$ such that for
$0<\mu\leq1$ and $N\geq K\mu^{-K}$, every collection of continuous complex functions
$f_0,f_1,f_2,f_3$ on $\R^3$ satisfies
\[
 \left\|A_N^{*j}(f_0,(f_i)_{i\ne j})
     -P_{(N^j)^{-1}(\mu^2)^{-1}}^{(j)}A_N^{*j}(f_0,(f_i)_{i\ne j})
          \right\|_{1/(1-\alpha_j)}
 \leq C\mu^c\|f_0\|_{1/\alpha_0}
                 \prod_{\substack{1\leq i\leq3\\i\ne j}}\|f_i\|_{1/\alpha_i}.
\]
The norms here are extended nonnegative norms, with precisely the zero-times-infinity
convention of \pid{ENNReal}. The unused input $f_j$ has no effect on the left or right side.
\end{corollary}
\begin{proof}
If any norm of an actually used input is zero, that continuous input is identically zero.
To justify this last statement, if $|f(x_0)|>0$, continuity bounds $|f|$ below by
$|f(x_0)|/2$ on a nonempty ball, whose Lebesgue measure is positive; its $L^p$ norm cannot
then be zero. Thus the adjoint is pointwise zero in this case, and so is its convolution.
Both sides are zero, even if another input norm is infinite.

Next suppose all used norms are nonzero. If any is infinite, their finite product is
infinite in $\mathbb R_{\geq0}\cup\{\infty\}$, and multiplication by the strictly positive
finite factor $C\mu^c$ keeps it infinite. The inequality is then immediate. It remains to
consider positive finite norms, where \cref{new:adjoint-all-exponents} applies.
For continuous inputs the $t$-integrand is continuous on the compact interval $[0,N]$ for
every $y$, so the concrete adjoint integral exists pointwise. The finite-norm estimate
identifies it almost everywhere with the $L^p$ output used there. Its one-coordinate
convolution is integrable for almost every output point by Young's inequality and Fubini;
at any exceptional points the totalized Bochner integral used by Lean does not affect the
extended norm. All functions involved are measurable (indeed the adjoint is continuous
locally by compact-parameter dominated convergence).

Finally $N>0$ and $\mu>0$ follow from the hypotheses, and real-power algebra gives
$(N^j)^{-1}(\mu^2)^{-1}=N^{-j}\mu^{-2}$. The mathematical output reciprocal is
$1-\alpha_j$, exactly as in the definition. Passing from a finite real $L^p$ norm to
\pid{eLpNorm} at \pid{ENNReal.ofReal p} preserves the estimate by its defining integral
formula. Multiplication, finite products, and positive real constants pass through
\pid{ENNReal.ofReal}; zero and infinity were already handled separately.
Replace mathematical indices $1,2,3$ by Lean indices $0,1,2$, and use
\pid{expo j = (j.val + 1)}. The product over the two other positive indices becomes
\pid{Finset.univ.erase j}. This is exactly the supplied definition, with no added
regularity, support, finite-norm, inverse-theorem, or improving-theorem hypothesis.
\end{proof}


\section{Implementation register and dependency audit}\label{new:register}
\subsection{What has and has not been checked}
This register refers to the submitted archive, not to a moving GitHub branch.
The toolchain is \pid{leanprover/lean4:v4.33.0-rc1}. The manifest pins mathlib to
\par{\small\texttt{79d0395a1825a6264ad5d269e35e60537518955e}}\par
and \pid{lean_spherical} to
\par{\small\texttt{d45bcdce0a6b68a3fea0581cba6dd84fdce31186}}.\par
The archive's SHA-256 is
\par{\small\texttt{f212ffa7214169734d3474228a44aec1f51a789eb14daae7a65b546f708c580b}}.\par

The source declarations listed below were read in the submitted files. This is source
inspection, not a fresh build or a transitive axiom check. No Lean executable or downloaded
pinned dependency tree was available in the working environment. The pinned Schwartz
Fourier source was independently accessible and checked; current online documentation is
used only to identify other library interfaces, not as a claim that all their argument lists
were checked at the pinned revision.

The document supplies mathematical proofs of every nonstandard analytic implication used
on this route. There are no residual assumptions named \pid{RealInverse},
\pid{HighestControl}, \pid{MajorArcProperty}, \pid{PETCertificate}, or \pid{KoszAdjoint}
in the final result. Turning the proof into Lean still requires elaborating the stated
measure, polynomial, and exponent interfaces. Names prefixed ``proposed'' below are an
implementation specification, not declarations claimed to have been added to the archive.

\subsection{The mathematical library leaves}
The following are the precise forms of the standard library facts at which the proof stops.
All specialized polynomial and harmonic-analysis consequences were established above.

\paragraph{Product integration and limits.}
For nonnegative measurable $F$ on two sigma-finite measure spaces, the product integral equals
both iterated integrals, including infinity. For integrable complex $F$, both iterated Bochner
integrals equal its product integral and its sections are integrable almost everywhere.
Dominated convergence applies to an almost-everywhere convergent sequence bounded in norm by
an integrable scalar function. Fatou applies to a sequence of nonnegative measurable functions.
The interfaces are in \pid{Mathlib.MeasureTheory.Integral.Prod},
\pid{Mathlib.MeasureTheory.Measure.Prod}, and the Bochner and Lebesgue integral modules.
The submitted file uses \pid{MeasureTheory.lintegral_prod},
\pid{MeasureTheory.lintegral_liminf_le},
\pid{MeasureTheory.StronglyMeasurable.integral_prod_right'}, and integrable Fubini.
All ``relative measures'' in the inverse section are normalized restrictions to its finite
boxes, not probability measures on all of $\R^n$.

\paragraph{Haar measure, regularity, and calculus.}
Translation preserves Lebesgue measure; an invertible linear change with matrix $A$ scales
it by $|\det A|$. Nonempty Euclidean balls have positive measure. For a finite-measure
measurable $E$ and $\epsilon>0$, there is compact $K\subseteq E$ with
$|E\setminus K|<\epsilon$. A compact $K$ has an open neighborhood $O$ with
$|O\setminus K|<\epsilon$. The regularity locators are
\pid{MeasurableSet.exists_isCompact_sdiff_lt} and
\pid{MeasurableSet.exists_isOpen_sdiff_lt} in
\pid{Mathlib.MeasureTheory.Measure.Regular}. Intersect $O$ with a bounded open box
containing $K$ for the bounded version used in \cref{new:translation-continuity}.
Approximation and translation continuity themselves were proved there.
The fundamental theorem of calculus, integration by parts, compactness, and maximum modulus
are the usual finite-dimensional calculus leaves. The supplied imports include
\pid{Mathlib.MeasureTheory.Integral.IntervalIntegral.IntegrationByParts} and
\pid{Mathlib.Analysis.Complex.Liouville}. Three lines and interpolation are proved in
\cref{new:three-lines,new:interpolation}, not imported as specialized analytic assertions.

\paragraph{Norms.}
H\"older, the integral triangle inequality, the Banach-space Bochner integral bound, and
completeness of $L^p$ for $1\leq p<\infty$ are the mean-inequalities, Bochner, and
\pid{LpSpace.Complete} interfaces already used by the archive. Their specialized forms are
in \cref{new:integral-inequalities,new:duality}. For $0<p<\infty$ the extended norm is
\[
 \operatorname{eLpNorm}(f,p)=\left(\int |f|^p\right)^{1/p}.
\]
The endpoint wrapper separates zero, finite positive, and infinite norms; it does not cancel
an extended norm or assume finiteness merely from continuity.

\paragraph{Schwartz inversion is the only Fourier inversion leaf.}
On Schwartz functions over a finite-dimensional real inner-product space, integral Fourier
transformation and inverse transformation preserve Schwartz space and are inverse linear maps.
Schwartz functions and their polynomially weighted derivatives are integrable. The checked
pinned module is \pid{Mathlib.Analysis.Distribution.SchwartzSpace.Fourier}, with
\pid{SchwartzMap.fourier_coe}, \pid{SchwartzMap.fourierInv_coe}, and its
\pid{FourierPair}/\pid{FourierInvPair} instances. The supplied
\pid{Auto.fourier_etaKer} uses \pid{FourierTransform.fourier_fourierInv_eq} on the Schwartz
cutoff. Plancherel for merely $L^1\cap L^2$ functions is \emph{not} a library leaf here:
\cref{new:integral-plancherel} proves it using the explicit approximate identity. An abstract
$L^2$ Fourier isometry and compatibility only for Schwartz functions would not suffice instead.

\paragraph{The exact Jacobian leaf.}
For a map $f:\R^n\to\R^n$ differentiable within a measurable set $S$ at every point of $S$
and injective on $S$, the one-sided inequality is
\[
 \int_S|\det Df(x)|\,dx\leq |f(S)|,
\]
with outer measure for the image when needed. The supplied source invokes
\pid{MeasureTheory.lintegral_abs_det_fderiv_le_addHaar_image} in
\pid{Auto.lintegral_abs_det_fderiv_le_card_mul_image}. In \cref{new:incidence-consequence}
there are exactly two explicitly described injective Borel pieces of a polynomial map.
Both images lie in the target set. No general polynomial multiplicity theorem or unproved
finite-fiber area formula is needed.

\paragraph{Polynomial and order algebra.}
The leaves are the binomial formula, coefficient equality, finite-product expansion,
independence of distinct monomials, the finite root bound for a nonzero scalar polynomial,
and well-founded induction on a fixed lexicographic tuple of natural numbers.
Use \pid{Polynomial} over \pid{MvPolynomial} with real coefficients, and carry the integer
coefficient certificate as an equality after coefficient mapping. Degree in the averaging
variable is distinct from total degree in the shift variables. All specialized consequences,
including the uniform termination bound, are proved in \cref{patch:pet-producer}.
The real sublevel estimate is proved in \cref{patch:multiaffine-sublevel}; no Remez theorem
is an additional premise.

\subsection{Existing declarations and their limits}
All line numbers below refer to
\pid{DFR/Auto/SmoothingIneq3D/Smoothing3D.lean} in the uploaded archive; they are locators,
not stable identifiers after editing.
\begin{longtable}{@{}p{0.52\linewidth}p{0.43\linewidth}@{}}
\toprule Declaration or group & Role and limitation \\
\midrule\endfirsthead
\toprule Declaration or group & Role and limitation \\
\midrule\endhead
\pid{Auto.E3}, \pid{Auto.expo} (122, 128) & Euclidean ambient space and natural exponent
\pid{j.val + 1}; \pid{expo} is not real-valued. \\
\pid{Auto.eta}, \pid{Auto.etaS}, \pid{Auto.etaKer};
\pid{Auto.fourier_etaKer} (795--858) & Chosen cutoff and inverse kernel. Do not replace it
silently by a monotone cutoff. \\
\pid{Auto.projKernel}, \pid{Auto.P}, \pid{Auto.fourier_P} (879--1186) & The existing Fourier
identity assumes continuity. The measurable generalization is proved in
\cref{new:multiplier-facts}. \\
\pid{Auto.P_self_adjoint}, \pid{Auto.Astar}, \pid{Auto.adjoint_identity} (1506--1619) & Match
all conjugations and the interval normalization exactly. \\
\pid{Auto.KoszAdjoint} (1803--1820) & Proposition, not proved theorem; target of
\cref{new:exact-lean-target}. \\
\pid{Auto.strong_real_improving'} (19225), \pid{Auto.kosz53_internal} (19430) & Existing
improving estimates; the instances needed here are proved in \cref{new:improving-section}. \\
\pid{Auto.locUnifPow_re_nonneg} (20097), \pid{Auto.cfgHeadInt_cube} (22515) & Positivity and
cube identities. PET must produce the actual configuration first. \\
\pid{Auto.locUnifPow_succ_split} (24596), \pid{Auto.locUnifPow_fdiff_eq} (24672) & Order and
difference expansions; retain all normalization ratios. \\
\pid{Auto.etaKer_decay} (26005), \pid{Auto.projKernel_decay} (26036),
\pid{Auto.integral_projKernel_mom} (26101), \pid{Auto.integral_projKernel_tail_le} (26119) &
Kernel moments and tails supporting \cref{new:kernel-facts}. \\
\pid{Auto.sq_re_gapScaled_le_locUnifPow} (37853) & Mixed-to-common scale comparison; the
polynomial state and sublevel selection are separate obligations. \\
\bottomrule
\end{longtable}

\subsection{New declarations in forward implementation order}
Each row names an interface whose full mathematical statement and proof occur earlier.
Rows may be split into additional elementary lemmas, but not replaced by axioms or proved
using later results that depend on them.
\begin{longtable}{@{}p{0.38\linewidth}p{0.57\linewidth}@{}}
\toprule Proposed interface & Proof and immediate dependencies \\
\midrule\endfirsthead
\toprule Proposed interface & Proof and immediate dependencies \\
\midrule\endhead
\pid{integralFourier_plancherel} & \Cref{new:translation-continuity,new:approximate-identity,new:integral-plancherel}:
regularity, Schwartz inversion, dominated convergence, Fatou, polarization. \\
\pid{directionalProjection_energy} & \Cref{new:multiplier-facts,new:average-multiplier}:
Fubini, preceding Plancherel, the real even cutoff. \\
\pid{multilinearInterpolation_simple} & \Cref{new:three-lines,new:interpolation}:
maximum modulus, explicit analytic simple-function families, duality. \\
\pid{realImproving_forAdjoint} & \Cref{new:flows,new:incidence-consequence,new:restricted-strong,new:improving}:
refinements, two flows, Jacobian inequality, dyadic summation, scaling. \\
\pid{signedIntegratedVdC}; \pid{cylinderCS} & \Cref{patch:signed-vdc,patch:cylinder}:
window-square expansion and Boolean-product induction. \\
\pid{protectedPet_update}; \pid{protectedPet_terminates} & \Cref{patch:pet-update,patch:pet-producer}:
symbolic state, leading-coefficient invariant, fixed pivot, lexicographic recursion. \\
\pid{affineHead_cube}; \pid{pet_gaps_nonzero} & \Cref{patch:cs-loss,patch:multiaffine-sublevel,patch:uniformize}:
ordered affine removal, signed losses, sublevel exclusion, positivity before comparison. \\
\pid{highestActiveControl} & \Cref{patch:highest-control}:
four phase differences, PET execution, and the separate degree-one case. \\
\pid{u2_selectFrequency}; \pid{dualDifference}; \pid{removeMissingPhases}; \pid{dummyPhase} &
\Cref{patch:u2-fourier-selection,patch:dual-difference,patch:missing-phase,patch:dummy-phase}:
Plancherel, least-rational selection, explicit cylinder inductions. \\
\pid{conditionalDegreeLowering} & \Cref{patch:conditional-degree}:
internal premise \pid{MajorArcProperty m l}; thresholds in the eight-step ledger. \\
\pid{realMajorArc}; \pid{structuredDegreeLowering} & \Cref{patch:major-arc,patch:structured-degree}:
induction on $m$ over every budget; only the smaller major-arc property is used. \\
\pid{removeHigherInputs}; \pid{lowestInputEnergy} & \Cref{patch:remove-high-inputs,patch:lowest-energy}:
backward elimination, useful frequency selection, scalar oscillation. \\
\pid{everyInputEnergy}; \pid{monomialInverse} & \Cref{patch:energy-core,new:monomial-inverse}:
lowest energy, freezing, shorter increasing degree lists, positive section energies. \\
\pid{compactAdjointHighPass} & \Cref{new:compact-highpass}:
inverse theorem applied to the truncated residual, multiplier annihilation, kernel tail. \\
\pid{compactAdjointDecay}; \pid{globalAdjointDecay} & \Cref{new:compact-decaying-point,new:global-decaying-point}:
interpolation with weight $60/61$, exact input-cube decomposition, summable off-diagonal sequence. \\
\pid{adjointAllExponents}; \pid{koszAdjoint} & \Cref{new:adjoint-all-exponents,new:exact-lean-target}:
explicit convex completion, approximation, zero/finite/infinite norm cases. \\
\bottomrule
\end{longtable}

\paragraph{The PET state is data, not an assumed certificate.}
Use a finite list of vector polynomials, a head index, and factor lists with records
(origin, conjugation bit, constant translation). Store proofs of normalized distinctness,
a singleton head list, maximal head degree, and the integer multilinear leading-coefficient
identity for every head gap. The update and preservation proofs are
\cref{patch:pet-update}. Store the type and length bound separately and recurse using $B(w,L)$.
Evaluate numerical shifts only after the formal execution. Good shifts then exist by the
proved sublevel and probability estimates. A structure containing the final cube conclusion
without these construction and preservation proofs would not implement the argument.

\paragraph{Analytic equality and extended norms.}
Develop the inverse module on bounded Borel functions and absolutely convergent integrals.
Use almost-everywhere equality when passing to $L^p$ classes. Keep \pid{Continuous} exactly
as required in the final wrapper; handle infinite norms before finite real-power algebra.
Do not assume that the indicator-truncated residual in \cref{new:compact-highpass} is
continuous: it generally is not, and its continuity is never used.

\paragraph{Final declaration and validation boundary.}
The proposed signature inside \pid{namespace Auto} is
\begin{verbatim}
theorem koszAdjoint (j : Fin 3) : KoszAdjoint j
\end{verbatim}
Its proof is \cref{new:exact-lean-target}, not an application of \pid{main_smoothing}, which
already depends on this proposition. After implementation, a build and a transitive axiom
inspection must check that the declaration has no added analytic axiom or \pid{sorryAx}.
Those checks have not been performed for this mathematical document. The existing subunit
estimate is separate and is not used to prove this declaration.

\subsection{Source correspondence and replaced steps}
\cite[Sections 4.2--4.5]{KoszEtAl} supplies the organization of the inverse argument:
uniformity control, conditional degree lowering, the major-arc induction, lowest-index
information, and general-index reduction. Lemma 4.43 omits detailed PET and points to the
method of \cite[Proposition 4.8]{KKL}. The real degree-at-most-three PET proof here is supplied
explicitly; no integer statement from that reference is asserted without proof to hold over
$\R$.

The improving argument follows the real refinement-and-flow method of
\cite[Section 5.6]{KoszEtAl}, with the two explicit flows and dyadic summation proved here.
The target is the real nontruncated monomial instance of \cite[Theorem 6.13]{KoszEtAl}.
The proof here replaces the structural Hahn--Banach route by \cref{new:compact-highpass}
and supplies support removal in \cref{new:global-decaying-point}. These are derivations in
this blueprint, not passages attributed verbatim to the paper.

The supplied \pid{patch_3_updated.tex} is the starting text for much of the inverse module.
The superseded unrestricted polynomial-to-cube assertion, generic degree-lowering assertion,
and circular use of the final inverse theorem do not occur on the proved dependency route.

\begin{thebibliography}{9}
\bibitem{KoszEtAl}
D. Kosz, M. Mirek, S. Peluse, R. Wan, and J. Wright,
\emph{The multilinear circle method and a question of Bergelson},
arXiv:2411.09478v4, 2 September 2026.
\url{https://arxiv.org/abs/2411.09478v4}.
\bibitem{KKL}
N. Kravitz, B. Kuca, and J. Leng,
\emph{Quantitative concatenation for polynomial box norms},
arXiv:2407.08636v2, especially Proposition 4.8.
\url{https://arxiv.org/abs/2407.08636v2}.
\bibitem{MathlibFourier}
The mathlib contributors, \emph{Schwartz-space Fourier transformation}, source at revision
\pid{79d0395a1825a6264ad5d269e35e60537518955e},
\pid{Mathlib/Analysis/Distribution/SchwartzSpace/Fourier.lean}.
\url{https://github.com/leanprover-community/mathlib4/blob/79d0395a1825a6264ad5d269e35e60537518955e/Mathlib/Analysis/Distribution/SchwartzSpace/Fourier.lean}.
\bibitem{MathlibRegular}
The mathlib contributors, \emph{Regular measures}, online API documentation,
\pid{Mathlib.MeasureTheory.Measure.Regular}. An interface locator, not a substitute for
checking the pinned dependency during implementation.
\url{https://leanprover-community.github.io/mathlib4_docs/Mathlib/MeasureTheory/Measure/Regular.html}.
\bibitem{SubmittedArchive}
\texttt{roos-j/lean-dfr}, submitted task-2 archive \pid{lean-dfr.task2.zip}, especially
\pid{DFR/Auto/SmoothingIneq3D/Smoothing3D.lean},
\pid{blueprints/patch_3_updated.tex}, the earlier blueprints, and automation ledgers.
Source version identified by the manifest and archive digest in \cref{new:register}.
\end{thebibliography}
\end{document}
